\documentclass[10pt,a4paper]{article}
\usepackage{fontspec,amsmath,amssymb,amsthm,mathtools,graphicx,bm,bbm,stmaryrd,nicefrac,mathrsfs,braket}
\usepackage[margin=22mm]{geometry}
\usepackage[bookmarks=false]{hyperref}
\setlength{\emergencystretch}{3em}
\title{Every nontrivial bipartition yields fast discrete colour refinement in random regular graphs}
\author{Mikhail Isaev, Tamás Makai, Brendan D. McKay, Paweł Prałat, Jane Tan, Maksim Zhukovskii}
\date{doi:10.4230/LIPIcs.ICALP.2026.114; CC BY 4.0}
\begin{document}
\maketitle
\noindent\textit{Editorial note.} Editable TeX converted from the official CC BY HTML prose and math alttext. Original official PDF/HTML retained. Only the specified claim counts; all author context and available source-local supporting proofs are included. Complete paper contains both dense and sparse proofs of the two supporting stages. The general main theorem is the sole collected claim.\par\medskip


\par\medskip\noindent\textbf{
1 Introduction}\par




Given an input graph $G$, a \emph{canonical labelling} algorithm computes a bijection $\pi_{G}:V(G)\mapsto\{1,\ldots,n\}$ with the following property: if a graph $G^{\prime}$ is isomorphic to $G$, then the relabelled versions of $G$ and $G^{\prime}$, under the actions of $\pi_{G}$ and $\pi_{G^{\prime}}$, are identical. There is a linear time reduction from the graph isomorphism problem to canonical labelling: once the labellings $\pi_{G},\pi_{G^{\prime}}$ have been computed for two input graphs $G,G^{\prime}$, it takes time $O(|E(G)|)$ to check whether $G,G^{\prime}$ are isomorphic. The best known (in the worst case) algorithm for the graph isomorphism problem is due to Babai [6, 27]: it runs in time $\exp(O(\log^{3}n))$ for $n$-vertex graphs. A quasi-polynomial bound $\exp(O(\log^{c}n))$ is also known for the canonical labelling problem [7]. Nevertheless, for graphs with bounded degree, both problems can be solved in polynomial time [9, 40]. In particular, this is the case for $d$-regular graphs when $d=\mathrm{const}.$ In this paper, we show that there exists a polynomial-time canonical labelling algorithm for \emph{almost all} $d$-regular graphs for \emph{all $0\leq d\leq n-1$}.\par




Colour refinement (CR) is a simple algorithmic routine that operates on vertex-coloured graphs. For an input graph
$G$ with initial colouring $C_{0}$, CR iteratively computes new colourings. At round $t$, $C_{t}(v)$ is a pair $(C_{t-1}(v),C_{t-1}(N(v)))$, where $C_{t-1}(N(v))$ is the multiset of $C_{t-1}$-colours of neighbours of $v$. That is, the process refines the initial partition $C_{0}$ and halts once the partition stabilises. Let us call a colouring \emph{discrete} if every pair of vertices is coloured differently. If CR runs on an uncoloured graph (i.e., there is only one initial colour) and outputs a discrete colouring, then, since the vertex colours are isomorphism-invariant, this yields a canonical labelling by numbering the colour names in the lexicographic order. In [8, 18, 25], it was proved that CR run on a binomial random graph $\mathcal{G}(n,p)$\footnote{The vertex set of $\mathcal{G}(n,p)$ is $\{1,\ldots,n\}$, and each pair of vertices is adjacent with probability $p=p(n)$, independently of the other pairs.} outputs a discrete colouring with high probability (\emph{whp}, in what follows)\footnote{A sequence of events $\mathcal{B}_{n}$ holds with high probability if $\mathbb{P}(\mathcal{B}_{n})\to 1$ as $n\to\infty$.} whenever $(1+\varepsilon)\frac{\ln n}{n}<p\leq\frac{1}{2}$, which implies a near linear time algorithm for canonical labelling of $G(n,p)$ [12].\par




We stress that for regular uncoloured graphs $G$, CR terminates immediately with a trivial colouring, and is therefore unsuitable for canonical labelling. For a positive integer $n$, we denote $[n]:=\{1,\ldots,n\}$. Let $d\leq n-1$ be a non-negative integer such that $dn$ is even. Let $\mathcal{G}_{n,d}$ be a uniform distribution over all $d$-regular graphs on $[n]$. We write $\mathbf{G}_{n}\sim\mathcal{G}_{n,d}$ for a graph sampled from this distribution, i.e. $\mathbf{G}_{n}$ is a uniformly random $d$-regular graph on $[n]$. Since, for constant $d$, efficient canonical labelling algorithms are known, we focus on the case $d=\omega(1)$. Moreover, since the edge complement of a $d$-regular graph is $(n-1-d)$-regular, the edge complement of $\mathcal{G}_{n,n-1-d}$ is distributed as $\mathcal{G}_{n,d}$. Therefore, we may restrict ourselves to $d\leq n/2$. Our main result shows that the triviality of the initial colouring is the only obstacle for a complete refinement: once a non-trivial initial colouring is produced, whp CR run on $\mathbf{G}_{n}\sim\mathcal{G}_{n,d}$ outputs a discrete colouring which is suitable for canonical labelling.\par




Before we state the main result of this paper, we need one more definition. For a connected graph $G$, we denote by $\operatorname{diam}(G)$ its diameter.\par




% SOURCE-ID Thmtheorem1


\par\medskip\noindent\textbf{
\textbf{Theorem 1}\textbf{.}
}\par



\emph{Let $d_{0}$ be large enough, let $d=d(n)$ be such that $d_{0}\leq d\leq n/2$, and let $\mathbf{G}_{n}\sim\mathcal{G}_{n,d}$. Then, the following holds whp:
for every non-trivial partition $[n]=V_{1}\sqcup V_{2}$ of the vertex set of $\mathbf{G}_{n}$, CR runs at most $2\operatorname{diam}(\mathbf{G}_{n})+3$ steps on $\mathbf{G}_{n}$ and outputs a discrete colouring.}\par



% END-SOURCE-ID Thmtheorem1




We note that we did not try to optimise the bound on the number of rounds, and we believe that $2\operatorname{diam}(\mathbf{G}_{n})+3$ is suboptimal. In particular, for $d\geq n^{1/2+\varepsilon}$, we show that $2\operatorname{diam}(\mathbf{G}_{n})+1=5$ rounds is enough. Actually, it is natural to suspect that the total number of rounds needed is $(1+o_{\operatorname{diam}(\mathbf{G}_{n})}(1))\operatorname{diam}(\mathbf{G}_{n})$ whp.\par




We now describe a possible approach to canonical labelling of $\mathbf{G}_{n}\sim\mathcal{G}_{n,d}$ based on our main result. Since the output of CR can be computed in time $O((n+|E(G)|)\log n)$ on an $n$-vertex graph $G$ [12], Theorem 1 reduces the problem to efficiently finding an isomorphism-preserving partition of $[n]$. To this end, recall that, for any $d=\omega(1)$, whp $\mathbf{G}_{n}$ contains a triangle [32]. Let $t_{i}$ be the number of triangles that contain the vertex $i\in[n]$ in $\mathbf{G}_{n}$, and let $t=\max\{t_{1},\ldots,t_{n}\}.$ Using fast square matrix multiplication, it is possible to compute the vector $(t_{1},\ldots,t_{n})$ in time $n^{\omega}$, where $\omega<2.372$ [1]. Alternatively, this vector can be computed in time $O(nd^{2})$ using the standard triangle-listing algorithm [16], which is faster when $d<n^{0.685}$, assuming the best known upper bound on the matrix multiplication exponent $\omega$ (and yields the bound $o(n^{2})$ for $d=o(\sqrt{n})$, which is faster than any algorithm based on square matrix multiplication since $\omega\geq 2$). Then, we can partition $[n]$ into sets $V_{1},V_{2}$ where $V_{1}$ contains all vertices $i$ for which $t_{i}=t$ and $V_{2}$ contains the rest.\par




Since whp $\mathbf{G}_{n}$ contains a triangle, $V_{1}$ is non-empty. When $d=o(n^{1/3})$, the number of triangles is sublinear whp (by Claim 9 below), which immediately implies that there are vertices that do not belong to a triangle, and so the set $V_{2}$ is non-empty as well. In order to show that $V_{2}$ is non-empty (i.e., that there are two vertices $u,v$ with $t_{u}\neq t_{v}$) whp when $d=\omega(n^{1/3})$, fix two vertices $u,v$ and expose their neighbourhoods $N(u),N(v)$. Then expose the edges inside $N(u)$. Assuming $t_{u}=t_{v}$, the set of exposed edges identifies the number of edges that have both endpoints in $N(v)$ but not in $N(u)$ (that is, in $E(\mathbf{G}_{n}[N(v)])\setminus E(\mathbf{G}_{n}[N(u)])$). By standard counting arguments (or using switchings), it follows directly that, for any fixed $x=x(n)\in\mathbb{Z}_{\geq 0}$, the probability that the latter set contains exactly $x$ edges tends to zero. Finally, when $d=\Theta(n^{1/3})$, it is known that the number of triangles $X_{n}$ in $\mathbf{G}_{n}$ satisfies a central limit theorem: $\frac{X_{n}-\mathbb{E}X_{n}}{\sqrt{\mathrm{Var}X_{n}}}$ converges in distribution to a standard normal random variable as $n\to\infty$ [21]. In particular, if $\mathbb{P}(\forall i\,\,t_{i}=t)>\varepsilon$, then $\mathbb{P}(3X_{n}/n\in\mathbb{Z})>\varepsilon$, contradicting the central limit theorem. Therefore, we get the following.\par



% SOURCE-ID Thmtheorem2


\par\medskip\noindent\textbf{
\textbf{Corollary 2}\textbf{.}
}\par



\emph{Let $d=d(n)$ be such that $d\to\infty$ and $n-d\to\infty$, and let $\mathbf{G}_{n}\sim\mathcal{G}_{n,d}$. There exists an algorithm that runs in time $O(\min\{n^{\omega},nd^{2}+nd\log n\})$ on $d$-regular $n$-vertex graphs and whp outputs a canonical labelling of $\mathbf{G}_{n}$.}\par



% END-SOURCE-ID Thmtheorem2



\par\medskip\noindent\textbf{
1.1 Related work}\par




In [13, 37], it is proved that, when $d\leq n^{\varepsilon},$ for a sufficiently small $\varepsilon$, then with high probability $\mathbf{G}_{n}\sim\mathcal{G}_{n,d}$ admits a canonical labelling via the 2-dimensional Weisfeiler-Leman algorithm (2-WL) [49], which is a generalisation of CR where the colouring is applied to pairs of vertices, and whose running time is $O(n^{3}\log n)$ [28]. The weaker version of 2-WL suggested by Bollobás [13]
canonically labels $\mathbf{G}_{n}$ in time $O(n^{3/2+\varepsilon})$ whp. Moreover, Kučera [37] claimed that his version of the algorithm runs in average time $O(nd)$ for $d=O(1)$ which we failed to verify\footnote{The paper does not provide a proof of this fact. However, the algorithm computes the vertices on shortest cycles as a subroutine and uses the following assertion [37, Theorem 3.1]: All cycles of the length $k$ in $d$-regular graph can be found in time $O(n\min\{n,(d-1)^{k/2}\})$. First, we believe that the factor $(d-1)^{k/2}$ should instead read $(d-1)^{\lceil k/2\rceil}$ (for instance, it is unclear how all triangles could be found in time $nd^{1.5}$, say in a union of $(d+1)$-cliques, there are $\Theta(nd^{2})$ triangles, which gives the lower bound $\Omega(nd^{2})$ on time needed to list them). Second, this bound (even with the fractional power) is not enough to get the expected time $O(nd)$. Indeed, $\mathbf{G}_{n}$ has a triangle with asymptotic probability $1-\exp(-(d-1)^{3}/6)$. So we get the bound on the expected time to be at least $(1-\exp(-(d-1)^{3}/6)-o_{n}(1))nd^{1.5}\sim nd^{1.5}$ as $d\to\infty$.}. We therefore conclude this paragraph by asking whether there exists a linear-time algorithm – running in time $O(nd)$ on $d$-regular $n$-vertex graphs – that canonically labels $\mathbf{G}_{n}$ whp, at least for some $3\leq d\leq n/2$.\par




For binomial random graphs $\mathcal{G}(n,p)$, the study of canonical labelling algorithms has been more extensive. Babai, Erdős, and Selkow [8] proved that CR outputs a canonical labelling of $\mathcal{G}(n,1/2)$ in linear time whp since it performs only a bounded number of refinement steps. The argument of [8] can be extended to show [14, Theorem 3.17] that the CR colouring of $\mathcal{G}(n,p)$ is whp discrete for all $n^{-1/5}\ln n\ll p\leq 1/2$. Bollobás [13] showed a polynomial time canonical labelling algorithm for $c_{1}\frac{\ln n}{n}\leq p\leq c_{2}\,n^{-11/12}$ for some positive constants $c_{1}$ and $c_{2}$, which is a weaker version of 2-WL. The next improvement was obtained by Czajka and Pandurangan [18]: they extended the range of applicability of CR to $\frac{\ln^{4}n}{n\ln\ln n}\ll p\leq\frac{1}{2}$, which was finally extended to $(1+\varepsilon)\frac{\ln n}{n}\leq p\leq\frac{1}{2}$ by Gaudio, Rácz, and Sridhar [25]. Linial and Mosheiff [39] showed that 2-WL outputs canonical labelling of $\mathcal{G}(n,p)$ whp when $\frac{1}{n}\ll p\leq\frac{1}{2}$. Finally, a polynomial time algorithm that labels canonically $\mathcal{G}(n,p)$ whp for all $0\leq p=p(n)\leq 1$ was independently established in [5, 47], with CR as the main ingredient.\par




A partition $[n]=V_{1}\sqcup\ldots\sqcup V_{t}$ of the vertex set of a graph $G$ is called \emph{equitable} if, for any $1\leq i,j\leq t$, any two vertices in $V_{i}$ have exactly the same number of neighbours in $V_{j}$. Clearly, if a graph $G$ admits a non-trivial equitable partition $V_{1}\sqcup\ldots\sqcup V_{t}$, then CR does not refine it. The opposite statement is also true – if there is a partition that CR does not refine, then this partition is equitable. Clearly, Theorem 1 implies that whp CR refines any non-trivial partition of $G_{n,d}$ with any number of parts. Therefore, it implies that whp $\mathbf{G}_{n}\sim\mathcal{G}_{n,d}$ does not have an equitable partition other than those with 1 part or $n$ parts. If a graph $G$ has a non-trivial automorphism group $\mathrm{Aut}(G)$, then for any non-trivial automorphism $\sigma\in\mathrm{Aut}(G)$, its cycle decomposition identifies an equitable partition of $G$.
Since no graph with three or more vertices has a
cyclic automorphism group acting without fixed points,
the absence of a non-trivial equitable partition implies
the absence of non-trivial automorphisms.
As a result, Theorem 1 implies that $\mathbf{G}_{n}$ is asymmetric whp for all $d$ such that $d\to\infty$ and $n-d\to\infty$. This is a known result for the entire range $3\leq n\leq d-4$, having been first established for $d=o(\sqrt{n})$ in [42] and then for $d\gg\log n$ in [31].\par





\par\medskip\noindent\textbf{
1.2 Proof strategy}\par




The key step to obtain Theorem 1 is to show that after some number of rounds of CR, one can coarsen the partition associated with colours to get a partition with $k$ parts of comparable size (for any desired arbitrarily large constant $k$). The following statement makes this precise.
Note that we may coarsen a partition at any stage of the CR algorithm if it is convenient for the argument that follows. Clearly, one may couple the original process with the modified one so that the partitions in the original process are refinements of the corresponding partitions in the modified one. In particular, if the modified process reaches discrete colouring, then so does the original one.\par



% SOURCE-ID Thmtheorem3


\par\medskip\noindent\textbf{
\textbf{Theorem 3}\textbf{.}
}\par



\emph{Let $d_{0}$ be large enough, let $d_{0}\leq d=d(n)\leq n/2$, and let $\mathbf{G}_{n}\sim\mathcal{G}_{n,d}$. Let $k\in\mathbb{N}$ be an arbitrary constant. Then, the following holds whp: for every non-trivial partition $[n]=V_{1}^{0}\sqcup V_{2}^{0}$ of the vertex set of $\mathbf{G}_{n}$, after $\operatorname{diam}(\mathbf{G}_{n})+2$ rounds of CR, there exists a partition $[n]=V_{1}\sqcup\ldots\sqcup V_{k}$ such that for any $i\in[k]$, $n/3k\leq|V_{i}|\leq 3n/k$ and $V_{i}$ is a union of some colour classes (in other words, $V_{1}\sqcup\ldots\sqcup V_{k}$ is a coarsening of the partition associated with resulting colour classes).}\par



% END-SOURCE-ID Thmtheorem3



Once there are many parts of linear and comparable sizes, CR distinguishes all vertices after some additional number of rounds.\par



% SOURCE-ID Thmtheorem4


\par\medskip\noindent\textbf{
\textbf{Theorem 4}\textbf{.}
}\par



\emph{Let $d_{0}$ be large enough, let $d_{0}\leq d=d(n)\leq n/2$, and let $\mathbf{G}_{n}\sim\mathcal{G}_{n,d}$. There exists some universal large constant $k\in\mathbb{N}$ such that the following holds whp: for every initial partition $[n]=V_{1}\sqcup\ldots\sqcup V_{k}$ such that for any $i\in[k]$, $n/3k\leq|V_{i}|\leq 3n/k$, CR terminates on $\mathbf{G}_{n}$ after at most $\operatorname{diam}(\mathbf{G}_{n})+1$ rounds and outputs a discrete colouring.}\par



% END-SOURCE-ID Thmtheorem4



Theorem 1 follows immediately from Theorem 3 and Theorem 4. To prove each of these two theorems, we address two regimes for $d$ separately: the dense case of $d\geq n^{1/2+\varepsilon}$, and the sparse case with $d_{0}\leq d\leq n^{10/17}$. The corresponding statements are reiterated in Sections 3, 4, 5, and 6.
Theorem 3 for $d\geq n^{1/2+\varepsilon}$ is proved in Section 3. Its proof consists of two parts. First, we show that whp after one refinement round there exists a coarsening $[n]=U_{1}\sqcup U_{2}$ of the CR-partition such that $|U_{1}|,|U_{2}|\gg n/d$, see Claim 16. One more round is needed to get all colour classes of size at most $\delta n$, for an arbitrary constant $\delta>0$, see Claim 17. The latter claim follows from the fact that there is no large set with all vertices having the same degree profile with respect to $(U_{1},U_{2})$. This is the main technical complication in the proof of Section 3 in the dense case: although this fact is easy to show in $\mathcal{G}(n,p)$, in random regular graphs we cannot rely on local limit theorems. Instead we use asymptotic estimations of the number of graphs with a given degree sequence as well as anti-concentration properties of hypergeometric distribution. The sparse case $d=o(n)$ is addressed in Section 4. Here, we show that, for every initial colouring $V_{1}\sqcup V_{2}$, where $|V_{1}|<cn$, for a sufficiently small constant $c>0$, after a few rounds of colour refinement, we will get a union of colour classes $U$ of size $|U|\in[cn/d,cn]$ (Lemma 14). One more round is needed to get a set $U^{\prime}$ of size $|U^{\prime}|\in[n\ell/d,0.999n]$, for an arbitrarily large constant $\ell$ (Lemma 22). Then, similarly to the dense case, we show that there is no set of size more than $\delta n$ such that all its vertices have same number of neighbours in $U^{\prime}$, that gives us the desired partition (Lemma 20). All three lemmas rely on switching arguments. In particular, the last two lemmas use switchings to establish an analogue of Erdős–Littlewood–Offord theorem in the context of uniformly random graphs with a fixed degree sequence.\par




Theorem 4 is proved in Sections 5 and 6. The dense case is significantly easier. For instance, when $d=\Theta(n)$, two refinement rounds are enough to obtain a discrete colouring whp. Indeed, let $u,v$ be two fixed vertices. Expose the neighbourhoods $N(u),N(v)$, and all the edges that touch $N(v)$. The exposed edges identify degree profiles of vertices in $N(u)\setminus N(v)$ with respect to the fixed partition. Since the latter set has size $\Theta(n)$, it is extremely unlikely that all the degrees are equal to the fixed values. Clearly, the probability of this event is $\left(1/\sqrt{n/k}\right)^{\Theta(kn)}$ in $\mathcal{G}(n,p)$, which is enough to overcome the union bound, with room to spare. In order to transfer this bound to random regular graphs, we use asymptotic enumeration of graphs with a given degree sequence and anti-concentration inequalities for hypergeometric distribution, as for the dense case in Theorem 3. For $n^{1/2+\varepsilon}\leq d=o(n)$, we need one additional refinement round in order to reach a set of vertices at distance at most 2 from $\{u,v\}$ of size $\Theta(n)$. The sparse case $d\leq n^{10/17}$, addressed in Section 6, requires a more delicate switching argument and constitutes the most technical part of the paper. Here, in order to reach a set of size $\Theta(n)$, from fixed vertices $u,v$, we need $\mathrm{diam}(\mathbf{G}_{n})$ rounds. Then, in contrast to the sparse case in Theorem 3, we need a multidimensional analogue of Erdős–Littlewood–Offord theorem (Claim 29), since the degree profiles are considered with respect to $k$ sets of the partition. Nevertheless, the claim can still be established by applying a similar switching argument $\Theta(k)$ times.\par





\par\medskip\noindent\textbf{
1.3 Organisation}\par




We start by presenting some preliminary results on properties of random graphs and concentration inequalities in Section 2. The rest of the paper is devoted to the proof of Theorems 3 and 4 which immediately imply our main Theorem 1. Theorem 3 is proved across Sections 3 and 4 where the dense and sparse cases are treated respectively. Sections 5 and 6 are devoted to the dense and sparse cases of Theorem 4.\par





\par\medskip\noindent\textbf{
1.4 Notation}\par




For a graph $G$, a set of vertices $U\subseteq V(G)$, and a non-negative integer $r$, we denote by $S_{r}(U)$ the sphere of radius $r$ around $U$ in the graph metric, omitting the dependency on $G$ since the underlying graph is always clear from the context. That is, $S_{r}(U)$ consists of vertices $v$ such that the length of a shortest path from $v$ to $U$ equals $r$. In particular, $S_{0}(U)=U$. We also denote $B_{r}(U)=\cup_{0\leq i\leq r}S_{i}(U)$ the ball of radius $r$ around $U$. We sometimes denote $S_{1}(U)$ by $N(U)$ and refer to it as the neighbourhood of $U$. For a set of vertices $X$ and a vertex $x\notin X$, we denote by $N_{X}(x)$ the number of neighbours of $x$ in $X$. We also use the standard notation $G[U]$ for the subgraph of $G$ induced by a set $U\subseteq V(G)$, and $G[U\times V]$ for the bipartite subgraph with (disjoint) parts $U$ and $V$, consisting of all edges of $G$ with one endpoint in $U$ and the other in $V$. We often write $A\sqcup B$ to denote the union of two \emph{disjoint} sets $A$ and $B$. Finally, for a random variable $X$ with distribution $Q$, we write $X\sim Q$. In particular, $X\sim\mathrm{Bin}(n,p)$ is a binomial random variable with $n$ trials and success probability $p$.\par






\par\medskip\noindent\textbf{
2 Preliminaries}\par




In this section, we collect some probabilistic tools as well as properties of random regular graphs that we will use in the main proofs to follow. Proofs of Lemma 5, Corollary 6, Lemma 8, and Lemma 14 appear in the extended version of the paper [29].\par




\par\medskip\noindent\textbf{
2.1 Concentration Inequalities}\par




We will use the following specific instances of Chernoff’s bound. Let $X\sim\textrm{Bin}(n,p)$. Then, a consequence of \emph{Chernoff’s bound} (see [30, Theorem 2.1]) is that for any $t\geq 0$ we have\par


\[\begin{gathered}
\displaystyle\mathbb{P}(X\geq\mathbb{E}X+t) \displaystyle\leq \displaystyle\exp\left(-\frac{t^{2}}{2(\mathbb{E}X+t/3)}\right) \\
\displaystyle\mathbb{P}(X\leq\mathbb{E}X-t) \displaystyle\leq \displaystyle\exp\left(-\frac{t^{2}}{2\mathbb{E}X}\right).
\end{gathered}\]


The same bounds hold for a random variable with the hypergeometric distribution with parameters $N$, $n$, and $m$ (see [30, Theorem 2.10]).\par





\par\medskip\noindent\textbf{
2.2 Properties of Binomials}\par




Let us start with a few auxiliary observations.\par



% SOURCE-ID Thmtheorem5


\par\medskip\noindent\textbf{
\textbf{Lemma 5}\textbf{.}
}\par



\emph{For all positive integers $b_{1}<a_{1}$, $b_{2}<a_{2}$,}\par


\[\begin{gathered}
\frac{a^{a}}{b^{b}(a-b)^{a-b}}\geq\frac{a_{1}^{a_{1}}a_{2}^{a_{2}}}{b_{1}^{b_{1}}(a_{1}-b_{1})^{a_{1}-b_{1}}b_{2}^{b_{2}}(a_{2}-b_{2})^{a_{2}-b_{2}}},
\end{gathered}\]


\emph{where $a=a_{1}+a_{2}$ and $b=b_{1}+b_{2}$.}\par



% END-SOURCE-ID Thmtheorem5



The preceding lemma has a useful corollary, which we record as follows.\par



% SOURCE-ID Thmtheorem6


\par\medskip\noindent\textbf{
\textbf{Corollary 6}\textbf{ }(Anti-concentration of hypergeometric distribution)\textbf{.}
}\par



\emph{For integers $0<b_{1}<a_{1}$, $0<b_{2}<a_{2}$,}\par


\[\begin{gathered}
\binom{a_{1}}{b_{1}}\binom{a_{2}}{b_{2}}\leq\binom{a}{b},
\end{gathered}\]


\emph{and}\par


\[\begin{gathered}
\binom{a_{1}}{b_{1}}\binom{a_{2}}{b_{2}}\leq\frac{2}{3}\sqrt{\frac{b(a-b)a_{1}a_{2}}{ab_{1}(a_{1}-b_{1})b_{2}(a_{2}-b_{2})}}\binom{a}{b},
\end{gathered}\]


\emph{where $a=a_{1}+a_{2}$ and $b=b_{1}+b_{2}$. In particular, for all integers $0<b<a$ and positive integers $k$,}\par


\[\begin{gathered}
\binom{a}{b}^{k}\leq\left(\frac{a}{b(a-b)}\right)^{(k-1)/2}\binom{ka}{kb}.
\end{gathered}\]


\emph{Moreover, for all integers $0<b<a$,}\par


\[\begin{gathered}
\binom{2a}{2b}\leq 4\sqrt{\frac{b(a-b)}{a}}\binom{a}{b}^{2}\leq 2\sqrt{a}\binom{a}{b}^{2}.
\end{gathered}\]



% END-SOURCE-ID Thmtheorem6




\par\medskip\noindent\textbf{
2.3 Counting Graphs}\par




For a given degree sequence $\mathbf{d}=(d_{1},\ldots,d_{n})$, we will use $g(\mathbf{d})$ to denote the number of graphs on the vertex set $[n]$ with the degree sequence $\mathbf{d}$. The following result gives precise (asymptotic) bounds on $g(\mathbf{d})$ (up to a multiplicative constant) for any degree sequence satisfying some mild condition. Dense graphs were investigated in [43] but the result was generalised to sparser graphs in [38].\par



% SOURCE-ID Thmtheorem7


\par\medskip\noindent\textbf{
\textbf{Theorem 7}\textbf{ }([38, 43, 44])\textbf{.}
}\par



\emph{Let $\mathbf{d}=(d_{1},\ldots,d_{n})$ be any degree sequence such that $\sum_{i=1}^{n}d_{i}$ is
even and for all $i\in[n]$, $|d_{i}-d|\leq d^{1/2+\varepsilon^{\prime}}$
for some $\varepsilon^{\prime}>0$, where $d$ is the average degree.
Let $m$ be the number of edges, and
$\eta=\frac{1}{n}\sum_{i=1}^{n}(d_{i}-d)^{2}$.
Suppose that $1\ll d\leq(1-\varepsilon_{0})n$ for some $\varepsilon_{0}>0$.
Then,}\par


\[\begin{gathered}
g(\mathbf{d})=\frac{\prod_{i=1}^{n}\binom{n-1}{d_{i}}}{m^{1/2}\binom{\binom{n}{2}}{m}}\exp\bigl(O(1)-\Theta(\eta^{2}/d^{2})\bigr).
\end{gathered}\]



% END-SOURCE-ID Thmtheorem7



We say that a sequence $\mathbf{d}=(d_{1},\ldots,d_{n})$ is \emph{balanced} if $|d_{i}-d_{j}|\leq 1$ for any $1\leq i<j\leq n$. The next observation is that $g(\mathbf{d})$ is maximized (over all sequences with a fixed even sum) when $\mathbf{d}$ is balanced.\par



% SOURCE-ID Thmtheorem8


\par\medskip\noindent\textbf{
\textbf{Lemma 8}\textbf{.}
}\par



\emph{Fix $m\in\left[\binom{n}{2}\right]$. The number of graphs on $n$ vertices and $m$ edges with a specified degree sequence $\mathbf{d}=(d_{1},\ldots,d_{n})$ (in particular, $\sum d_{i}=2m$) is maximized when the degree sequence is as even as possible. In other words,}\par


\[\begin{gathered}
\max\left\{g(\mathbf{d}):\sum d_{i}=2m\right\}=g(\mathbf{\hat{d}}),
\end{gathered}\]


\emph{where $\mathbf{\hat{d}}$ is the degree sequence, unique up to order, with only $\lfloor\sum d_{i}/n\rfloor$ and $\lceil\sum d_{i}/n\rceil$.}\par



% END-SOURCE-ID Thmtheorem8



We also recall the probability bound on the event that a uniformly random graph with a given degree sequence contains a specified set of edges.\par



% SOURCE-ID Thmtheorem9


\par\medskip\noindent\textbf{
Claim 9 ([41]). 
}\par



Let $\mathbf{G}_{n}$ be a uniformly random graph on the vertex set $[n]$ with a fixed degree sequence $(d_{1},\ldots,d_{n})$. Let $H$ be a graph on $[n]$ with degree sequence $(d^{\prime}_{1},\ldots,d^{\prime}_{n})$ such that $d^{\prime}_{i}\leq d_{i}$ for all $i\in[n]$. Let $m=\frac{1}{2}\sum d_{i}$ and $m^{\prime}=\frac{1}{2}\sum d^{\prime}_{i}$ be the number of edges in $\mathbf{G}_{n}$ and $H$, respectively. Let $d=\max\{d_{1},\ldots,d_{n}\}=o(m^{1/2})$ and let $m^{\prime}\leq m/2$. Then,\par


\[\begin{gathered}
\mathbb{P}(H\subseteq\mathbf{G}_{n})\leq\frac{\prod_{i=1}^{n}d_{i}(d_{i}-1)\ldots(d_{i}-d^{\prime}_{i}+1)}{2^{m^{\prime}}(m-2d^{2})(m-2d^{2}-1)\ldots(m-2d^{2}-m^{\prime}+1)}.
\end{gathered}\]



% END-SOURCE-ID Thmtheorem9



We note that Claim 9 immediately implies the following.\par



% SOURCE-ID Thmtheorem10


\par\medskip\noindent\textbf{
Claim 10. 
}\par



Under the assumptions of Claim 9, for all $n$ large enough when the degree sequence is regular,\par


\[\begin{gathered}
\mathbb{P}(H\subseteq\mathbf{G}_{n})\leq\left(2d/n\right)^{|E(H)|}.
\end{gathered}\]



% END-SOURCE-ID Thmtheorem10




\par\medskip\noindent\textbf{
2.4 Sandwiching Graphs}\par




Consider the binomial random graph $\mathcal{G}(n,p)$ which has vertex set $[n]$ and each potential edge is included independently at random with probability $p$; $p=p(n)$ could be, and usually is, a function of $n$ that tends to zero as $n\to\infty$. Since the independence of the edges allows the use of a wide variety of techniques, $\mathcal{G}(n,p)$ is typically much easier to study compared to $\mathcal{G}_{n,d}$. As a result, it is tempting to hope for a general purpose “black box” theorem that is able to translate results between $\mathcal{G}(n,p)$ and $\mathcal{G}_{n,d}$. In 2004, Kim and Vu [33] formalized this desire in their famous “sandwich conjecture”. After more than 20 years and a number of important contributions [19, 22, 23, 34, 24], the conjecture was finally proved [11].\par



% SOURCE-ID Thmtheorem11


\par\medskip\noindent\textbf{
\textbf{Theorem 11}\textbf{ }(Theorem 1.1 [11])\textbf{.}
}\par



\emph{For each $\epsilon>0$ there is some $C>0$ such that the following holds for each $d\geq C\log n$. There is a coupling $(\mathbf{G}_{*},\mathbf{G},\mathbf{G}^{*})$ of random graphs such that $\mathbf{G}_{*}\sim\mathcal{G}(n,(1-\epsilon)d/n)$, $\mathbf{G}\sim\mathcal{G}_{n,d}$, $\mathbf{G}^{*}\sim\mathcal{G}(n,(1+\epsilon)d/n)$, and whp $\mathbf{G}_{*}\subset\mathbf{G}\subset\mathbf{G}^{*}$.}\par



% END-SOURCE-ID Thmtheorem11




\par\medskip\noindent\textbf{
2.5 Expansion Properties}\par




We will use the expansion properties of random $d$-regular graphs that follow from their eigenvalues. The adjacency matrix $A=A(G)$
of a given $d$-regular graph $G$ on $n$ vertices, is an $n\times n$ real symmetric matrix. Thus, the matrix $A$ has $n$ real eigenvalues which we denote by $d=\lambda_{1}\geq\lambda_{2}\geq\cdots\geq\lambda_{n}$. It is known that several structural properties of a $d$-regular graph are reflected in its spectrum. Since we focus on expansion properties, we are particularly interested in the following quantity: $\lambda=\lambda(G):=\max\{|\lambda_{2}|,|\lambda_{n}|\}$.\par




The number of edges $e(A,B)$ between two sets $A$ and $B$ in a random $d$-regular graph on $n$ vertices is expected to be close to $d|A||B|/n$. (Note that $A\cap B$ does not have to be empty; in general, $e(A,B)$ is defined to be the number of edges between $A\setminus B$ to $B$ plus twice the number of edges that contain only vertices of $A\cap B$.) A small $\lambda$ (that is, a large spectral gap) implies that the deviation is small. The following bound is very convenient.\par



% SOURCE-ID Thmtheorem12


\par\medskip\noindent\textbf{
\textbf{Theorem 12}\textbf{ }([3, 35])\textbf{.}
}\par



\emph{Let $G$ be a $d$-regular graph. Then for any two sets of vertices $A,B\subseteq V(G)$, the number $e(A,B)$ of edges of $G$ with one endpoint in $A$ and another endpoint in $B$ satisfies}\par


\[\begin{gathered}
\left|e(A,B)-\frac{d|A||B|}{n}\right|\leq\lambda\sqrt{|A||B|}.
\end{gathered}\]



% END-SOURCE-ID Thmtheorem12



We will apply the Expander Mixing Lemma (Theorem 12) together with an asymptotic bound on $\lambda$ for random regular graphs. It was first established by Friedman [20] for constant $d\geq 3$, confirming the conjecture of Alon [2]. The case of $d\to\infty$ was then conjectured by Vu [48]. After a series of important contributions [4, 15, 17, 36, 46], it was resolved for all $d=o(n)$ by Bauerschmidt, Huang, Knowles, and Yau [10] and Sarid [45], and then for $d=\Theta(n)$ by He [26].\par



% SOURCE-ID Thmtheorem13


\par\medskip\noindent\textbf{
\textbf{Theorem 13}\textbf{ }([10, 26, 45])\textbf{.}
}\par



\emph{Let $3\leq d\leq n/2$ and $\mathbf{G}_{n}\sim\mathcal{G}_{n,d}$. Then, whp $\lambda(\mathbf{G}_{n})\leq(2+o(1))\sqrt{d(1-d/n)}$.}\par



% END-SOURCE-ID Thmtheorem13



We will also require a finer expansion result ensuring that, for every set, its size remains concentrated after several rounds of expansion.\par



% SOURCE-ID Thmtheorem14


\par\medskip\noindent\textbf{
\textbf{Lemma 14}\textbf{.}
}\par



\emph{Let $c>0$ be small enough and $d_{0}\in\mathbb{N}$ be large enough (independent of $c$). Let $d_{0}\leq d\leq n/2$ and $\mathbf{G}_{n}\sim\mathcal{G}_{n,d}$. Then, the following holds whp:
for every set $U$ of size $u=|U|\leq\frac{cn}{d}$ and for every positive integer $r$ such that $ud(d-1)^{r-1}\leq cn$,}\par


\[\begin{gathered}
|S_{r}(U)|\geq(1-100c-4\ln d/d)ud(d-1)^{r-1}.
\end{gathered}\]



% END-SOURCE-ID Thmtheorem14





\par\medskip\noindent\textbf{
3 Proof of Theorem 3: Dense Case}\par




Here we prove the following.\par



% SOURCE-ID Thmtheorem15


\par\medskip\noindent\textbf{
\textbf{Theorem 15}\textbf{.}
}\par



\emph{Let $\varepsilon\in(0,1/2)$, $n^{1/2+\varepsilon}\leq d=d(n)\leq n/2$, and let $\mathbf{G}_{n}\sim\mathcal{G}_{n,d}$. Let $k\in\mathbb{N}$ be an arbitrary constant. Then, the following holds whp: for every non-trivial partition $[n]=V_{1}^{0}\sqcup V_{2}^{0}$ of the vertex set of $\mathbf{G}_{n}$, after two rounds of CR, there exists a partition $[n]=V_{1}\sqcup\ldots\sqcup V_{k}$ such that for any $i\in[k]$, $n/3k\leq|V_{i}|\leq 3n/k$ and $V_{i}$ is a union of some colour classes.}\par



% END-SOURCE-ID Thmtheorem15



Theorem 15 follows easily from the following two claims.\par



% SOURCE-ID Thmtheorem16


\par\medskip\noindent\textbf{
Claim 16. 
}\par



Let $\varepsilon\in(0,1/2)$, $n^{1/2+\varepsilon}\leq d=d(n)\leq n/2$, and let $\mathbf{G}_{n}\sim\mathcal{G}_{n,d}$.
For every $C>1$, the following property holds whp: for every non-trivial partition $[n]=V_{1}\sqcup V_{2}$ with $\min\{|V_{1}|,|V_{2}|\}<Cn/d$, after one round of CR, there exists a partition $[n]=U_{1}\sqcup U_{2}$ with $\min\{|U_{1}|,|U_{2}|\}\geq Cn/d$ such that $U_{1}$ and $U_{2}$ are unions of some colour classes (in other words, $U_{1}\sqcup U_{2}$ is a coarsening of the partition
associated with resulting colour classes).\par



% END-SOURCE-ID Thmtheorem16


% SOURCE-ID Thmtheorem17


\par\medskip\noindent\textbf{
Claim 17. 
}\par



Let $\varepsilon\in(0,1/2)$, $n^{1/2+\varepsilon}\leq d=d(n)\leq n/2$, and let $\mathbf{G}_{n}\sim\mathcal{G}_{n,d}$.
For any $\delta\in(0,1]$, there exists $C=C(\delta)>0$ such that the following property holds whp: for every non-trivial partition $[n]=V_{1}\sqcup V_{2}$ with $\min\{|V_{1}|,|V_{2}|\}\geq Cn/d$, there is no colour class of size more than $\delta n$ after one round of CR.\par



% END-SOURCE-ID Thmtheorem17



The proof of the first claim is fairly straightforward and relies on the “sandwich theorem” (Theorem 11); it appears in the extended version of the paper [29]. Before we prove the second claim, let us show how they imply Theorem 15.\par




\par\medskip\noindent\textbf{Proof of Theorem 15.}\par




Since we aim for the statement that holds whp, we may assume that the statements in Claim 16 and in Claim 17 hold deterministically. Fix any $k\in\mathbb{N}$, and let $\delta=1/3k$. Let $C=C(\delta)$ be the large enough constant implied by Claim 17.\par




Consider any non-trivial partition $[n]=V_{1}^{0}\sqcup V_{2}^{0}$. If $\min\{|V_{1}^{0}|,|V_{2}^{0}|\}<Cn/d$, then after one round of CR (and coarsening), we get a partition into two colour classes where both of the colour classes have size at least $Cn/d$ (by Claim 16). After another round of CR, all colour classes have size at most $\delta n=n/3k$ (by Claim 17). If $\min\{|V_{1}^{0}|,|V_{2}^{0}|\}\geq Cn/d$, then we get the above property after a single round of CR.\par




To get the desired partition into $k$ parts, each of size at least $n/3k$, one can iteratively merge any two colour classes of size at most $\delta n$ until there is at most one such class remaining. After possibly merging this last class (if it exists) with any other arbitrarily chosen class, we get at least $k=1/3\delta$ classes (but at most $3k$ of them), each of size at least $\delta n=n/3k$ but at most $3\delta n$. Finally, if there are more than $k$ classes, one can arbitrarily merge some triples of them (and, perhaps, one pair) to get exactly $k$ classes, each of size at most $9\delta n=3n/k$. This finishes the proof of the theorem.
$\blacktriangleleft$\par




In order to prove Claim 17, we will make use of the following simple observation (see its proof in the extended version of the paper [29]).\par



% SOURCE-ID Thmtheorem18


\par\medskip\noindent\textbf{
\textbf{Lemma 18}\textbf{.}
}\par



\emph{Let $\varepsilon\in(0,1/2)$, $n^{1/2+\varepsilon}\leq d=d(n)\leq n/2$, and let $\mathbf{G}_{n}\sim\mathcal{G}_{n,d}$.
For any $\delta>0$, the following property holds whp: for any $U\subseteq[n]$ of size $u=|U|>\delta n$ and any $V\subseteq[n]\setminus U$ of size $v=|V|\geq 500n/(d\delta)$, the number of edges $e(U,V)$ between $U$ and $V$ satisfies the following bounds}\par


\[\begin{gathered}
0.8uv\cdot\frac{d}{n-1}\leq e(U,V)\leq 1.2uv\cdot\frac{d}{n-1}\,.
\end{gathered}\]



% END-SOURCE-ID Thmtheorem18




\par\medskip\noindent\textbf{
\textbf{Proof of Claim 17}.}\par




Fix any $\delta\in(0,1]$ and let $C=C(\delta)$ be a large enough constant that will be specified later. In particular, we will assume that $C\geq 500/\delta$ so that we may apply Lemma 18.\par




Suppose that there exists a partition $[n]=V\sqcup([n]\setminus V)$ with $Cn/d\leq|V|\leq n/2$ such that after one round of CR there exists a colour class $U$ of size more than $\delta n$. Note that this implies that every vertex in $U$ has the same number of neighbours in $V$ (hence every vertex in $U$ also has the same number of neighbours in $[n]\setminus V$). If $U\subseteq V$, then it will be convenient to concentrate on the number of neighbours in $[n]\setminus V\subseteq[n]\setminus U$ but if $U\subseteq[n]\setminus V$, then we will concentrate on the number of neighbours in $V\subseteq[n]\setminus U$. Our goal is to estimate the probability of the weaker but necessary property that there exists a pair of sets $(U,V)$ such that $V\subseteq[n]\setminus U$, $|V|\geq Cn/d$, $|U|>\delta n$, and every vertex in $U$ has the same number of neighbours in $V$.\par




Fix $V\subseteq[n]$ and $U\subseteq[n]\setminus V$ such that $v:=|V|\geq Cn/d$ and $u:=|U|>\delta n$. Note that, in particular, $v<(1-\delta)n$. For each non-negative integer $k\leq v$, define $\mathcal{E}_{k}(U,V)$ to be the event that every vertex in $U$ has exactly $k$ neighbours in $V$. By Lemma 18, since we aim for a statement that holds whp, we may assume that the number of edges between $U$ and $V$ is at least $0.8uvd/(n-1)$ and at most $1.2uvd/(n-1)$. Hence, we may restrict to considering $k$ such that\par


\[\begin{gathered}
0.8v\cdot\frac{d}{n-1}\nobreak\ \leq\nobreak\ k\nobreak\ \leq\nobreak\ 1.2v\cdot\frac{d}{n-1}.
\end{gathered}\]




First, note that the expected number of edges induced by $[n]\setminus V$ is $\binom{n-v}{2}\cdot\frac{d}{n-1}$. We will show that it is highly unlikely that the actual number deviates substantially from it. Let\par


\[\begin{gathered}
m_{-}=0.9\binom{n-v}{2}\cdot\frac{d}{n-1}\qquad\text{ and }\qquad m_{+}=1.1\binom{n-v}{2}\cdot\frac{d}{n-1}.
\end{gathered}\]


By Theorem 7 and the Stirling’s formula ($s!=(1+o(1))\sqrt{2\pi s}(s/e)^{s}$), letting $\mathbf{d}:=(d,\ldots,d)\in\mathbb{Z}^{n}$, the number of $d$-regular graphs on $[n]$ can be estimated as follows:\par


\[\begin{gathered}
\displaystyle g(\mathbf{d}) \displaystyle=\Theta\left(\frac{{n-1\choose d}^{n}}{\sqrt{dn}{{\binom{n}{2}}\choose dn/2}}\right) \\
\displaystyle=\left(\!\sqrt{\frac{n}{2\pi d(n-d)}}\cdot\frac{(1+o(1))\left(\frac{n-1}{e}\right)^{n-1}}{\left(\frac{d}{e}\right)^{d}\left(\frac{n-1-d}{e}\right)^{n-1-d}}\right)^{n}\!\left(\!\frac{\left(\frac{dn/2}{e}\right)^{dn/2}\left(\frac{n(n-1-d)/2}{e}\right)^{n(n-1-d)/2}}{\Theta(1)\left(\frac{n(n-1)/2}{e}\right)^{n(n-1)/2}}\right) \\
\displaystyle=\Theta(1)\cdot{{\binom{n}{2}}\choose dn/2}\left(\!(1+o(1))\sqrt{\frac{n}{2\pi d(n-d)}}\right)^{n}={{\binom{n}{2}}\choose dn/2}d^{-\Theta(n)}.
\end{gathered}\]


Hence, the probability that the number of edges induced by $[n]\setminus V$ is at most $m_{-}$ or at least $m_{+}$ can be upper bounded by\par


\[\begin{gathered}
\displaystyle\sum_{\begin{subarray}{c}m\leq m_{-}\\
m\geq m_{+}\end{subarray}}\frac{{{n-v\choose 2}\choose m}{{\binom{n}{2}}-{n-v\choose 2}\choose dn/2-m}}{g(\mathbf{d})} \displaystyle=n^{\Theta(n)}\sum_{\begin{subarray}{c}m\leq m_{-}\\
m\geq m_{+}\end{subarray}}\frac{{{n-v\choose 2}\choose m}{{\binom{n}{2}}-{n-v\choose 2}\choose dn/2-m}}{{{\binom{n}{2}}\choose dn/2}} \\
\displaystyle=n^{\Theta(n)}\cdot\mathbb{P}\left(\eta\leq m_{-}\text{ or }\eta\geq m_{+}\right),
\end{gathered}\]


where $\eta$ is the hypergeometric random variable with parameters $\binom{n}{2}$, $\binom{n-v}{2}$, and $dn/2$. Clearly,\par


\[\begin{gathered}
\mathbb{E}\eta=\frac{dn}{2}\cdot\frac{\binom{n-v}{2}}{\binom{n}{2}}={n-v\choose 2}\frac{d}{n-1}=\Theta(dn).
\end{gathered}\]


By Chernoff’s bound for hypergeometric distribution (see the comment right after (1), (2)),\par


\[\begin{gathered}
\mathbb{P}(\eta\leq m_{-}\text{ or }\eta\geq m_{+})=\mathbb{P}\left(|\eta-\mathbb{E}\eta|\geq 0.1\mathbb{E}\eta\right)=\exp\left(-\Omega\left(\mathbb{E}\eta\right)\right)=\exp(-\Omega(dn)).
\end{gathered}\]




Similarly, if $|V|\geq n^{3/4}$, then the expected number of edges induced by $V$ is $\binom{v}{2}\cdot\frac{d}{n-1}=\Theta(v^{2}d/n)$ and we get that with probability $\exp(-\Omega(dv^{2}/n))=\exp(-\Omega(dn^{1/2}))$, the number of edges induced by $V$ is at most $0.9\binom{v}{2}\cdot\frac{d}{n-1}$ or at least $1.1\binom{v}{2}\cdot\frac{d}{n-1}$.\par




It remains to concentrate on the case when the number of edges induced by $[n]\setminus V$ is between $m_{-}$ and $m_{+}$, that is, when the average degree of the graph induced by $[n]\setminus V$ is at least $0.9(n-1-v)\cdot\frac{d}{n-1}$ but at most $1.1(n-1-v)\cdot\frac{d}{n-1}$. Let us first deal with the case when $|V|\geq n^{3/4}$ so we may additionally assume that the average degree of the graph induced by $V$ is at least $0.9(v-1)\cdot\frac{d}{n-1}$ but at most $1.1(v-1)\cdot\frac{d}{n-1}$. By Theorem 7 and Lemma 8,\par


\[\begin{gathered}
\mathbb{P}(\mathcal{E}_{k}(U,V))=e^{-\Omega(dn)}+e^{-\Omega(dn^{1/2})}+O\left(\sum_{D}h(k,u,v,D)\right),
\end{gathered}\]


where $D$ denotes the number of edges between $V$ and $[n]\setminus(V\cup U)$, and\par


\[\begin{gathered}
h(k,u,v,D):=\frac{{v\choose k}^{u}{(n-v-u)v\choose D}\binom{v-1}{d-(D+ku)/v}^{v}{n-1-v\choose d-(D+ku)/(n-v)}^{n-v}}{\binom{{v\choose 2}}{(dv-ku-D)/2}{{n-v\choose 2}\choose(d(n-v)-ku-D)/2}g(\mathbf{d})}\,.
\end{gathered}\]


Indeed, there are at most ${v\choose k}^{u}$ ways to place edges between $U$ and $V$, and at most ${(n-v-u)v\choose D}$ ways to place edges between $V$ and $[n]\setminus(V\cup U)$. (Note that these values are trivial upper bounds but not the exact ones as some choices create vertices of degree more than $d$.) It remains to estimate the number of graphs induced by the set $[n]\setminus V$ and the number of graphs induced by the set $V$. Importantly, once other edges are fixed, these graphs have a fixed degree distribution. In particular, the average degree of the graphs induced by $[n]\setminus V$ is precisely $f(D):=d-(D+ku)/(n-v)$. Similarly, the average degree of the graphs induced by $V$ is $\tilde{f}(D):=d-(D+ku)/v$. Hence, we may use Theorem 7 and Lemma 8 to get upper bounds for the number of such graphs. (Let us point out that $f(D)$ and $\tilde{f}(D)$ are not necessarily integers. However, to keep the notation simple, we write ${n-1-v\choose f(D)}^{n-v}$ instead of the product of $n-v$ terms, each of them being ${n-1-v\choose\lfloor f(D)\rfloor}$ or ${n-1-v\choose\lceil f(D)\rceil}$.) Finally, since the average degree of the graph induced by $[n]\setminus V$ and the one induced by $V$ are restricted, $D$ satisfies the requirements\par


\[\begin{gathered}
\displaystyle 0.9(n-1-v)\cdot\frac{d}{n-1}\leq \displaystyle f(D) \displaystyle\leq 1.1(n-1-v)\cdot\frac{d}{n-1} \\
\displaystyle 0.9(v-1)\cdot\frac{d}{n-1}\leq \displaystyle\tilde{f}(D) \displaystyle\leq 1.1(v-1)\cdot\frac{d}{n-1}.
\end{gathered}\]




There are three binomials in the numerator of $h(k,u,v,D)$ that are raised to powers that are functions of $n$. We need to take advantage of them using Corollary 6 (see (5)). By (9),\par


\[\begin{gathered}
\displaystyle{n-1-v\choose f(D)}^{n-v} \displaystyle\leq \displaystyle\left(\frac{(1+o(1))(n-1-v)}{f(D)(n-1-v-f(D))}\right)^{(n-v-1)/2}\binom{(n-v)(n-v-1)}{d(n-v)-ku-D} \\
\displaystyle\leq \displaystyle\left(\frac{1+o(1)}{0.9\cdot\delta\cdot d\cdot(1-1.1\cdot 0.5)}\right)^{(n-v-1)/2}\binom{(n-v)(n-v-1)}{d(n-v)-ku-D} \\
\displaystyle\leq \displaystyle\left(\frac{3}{\delta d}\right)^{(n-v-1)/2}\binom{(n-v)(n-v-1)}{d(n-v)-ku-D}.
\end{gathered}\]


Similarly, by (10),\par


\[\begin{gathered}
\displaystyle\binom{v-1}{\tilde{f}(D)}^{v} \displaystyle\leq \displaystyle\left(\frac{(1+o(1))(v-1)}{\tilde{f}(D)(v-1-\tilde{f}(D))}\right)^{(v-1)/2}\binom{v(v-1)}{dv-ku-D} \\
\displaystyle\leq \displaystyle\left(\frac{1+o(1)}{0.9\cdot(vd/n)\cdot(1-1.1\cdot 0.5)}\right)^{(v-1)/2}\binom{v(v-1)}{dv-ku-D} \\
\displaystyle\leq \displaystyle\left(\frac{3n}{vd}\right)^{(v-1)/2}\binom{v(v-1)}{dv-ku-D}.
\end{gathered}\]


Finally, by (7),\par


\[\begin{gathered}
\displaystyle\binom{v}{k}^{u} \displaystyle\leq \displaystyle\left(\frac{v}{k(v-k)}\right)^{(u-1)/2}\binom{vu}{ku}\leq\left(\frac{1+o(1)}{0.8vd/n(1-1.2\cdot 0.5)}\right)^{(u-1)/2}\binom{vu}{ku} \\
\displaystyle\leq \displaystyle\left(\frac{4n}{dv}\right)^{(u-1)/2}\binom{vu}{ku}.
\end{gathered}\]




For future reference, let us highlight that (12) only holds when $v\geq n^{3/4}$, whilst the other two bounds (11) and (13) hold in general. Substituting in these three bounds, we get\par


\[\begin{gathered}
h(k,u,v,D)=O\left(n^{\frac{3}{2}}\ \frac{\left(\frac{4n}{dv}\right)^{\frac{u}{2}}\binom{vu}{ku}{(n-v-u)v\choose D}\left(\frac{3n}{vd}\right)^{\frac{v}{2}}\binom{v(v-1)}{dv-ku-D}\left(\frac{3}{\delta d}\right)^{\frac{n-v}{2}}\binom{(n-v)(n-v-1)}{d(n-v)-ku-D}}{\binom{{v\choose 2}}{(dv-ku-D)/2}{{n-v\choose 2}\choose(d(n-v)-ku-D)/2}g(\mathbf{d})}\right).
\end{gathered}\]


Now, by Corollary 6 (see (6)), the latter quantity equals\par


\[\begin{gathered}
O\left(n^{7/2}\ \frac{\left(\frac{4n}{dv}\right)^{u/2}\binom{vu}{ku}{(n-v-u)v\choose D}\left(\frac{3n}{vd}\right)^{v/2}\binom{{v\choose 2}}{(dv-ku-D)/2}\left(\frac{3}{\delta d}\right)^{(n-v)/2}{{n-v\choose 2}\choose(d(n-v)-ku-D)/2}}{g(\mathbf{d})}\right).
\end{gathered}\]


By Corollary 6 (see (3)), we can collect all binomial coefficients together to get\par


\[\begin{gathered}
\displaystyle h(k,u,v,D) \displaystyle= \displaystyle O\left(n^{\frac{7}{2}}\ \frac{\left(\frac{4n}{dv}\right)^{\frac{u}{2}}\binom{vu+(n-v-u)v+{v\choose 2}+{n-v\choose 2}}{ku+D+(dv-ku-D)/2+(d(n-v)-ku-D)/2}\left(\frac{3n}{vd}\right)^{\frac{v}{2}}\left(\frac{3}{\delta d}\right)^{\frac{n-v}{2}}}{g(\mathbf{d})}\right) \\
\displaystyle= \displaystyle O\left(n^{7/2}\ \frac{\left(\frac{4n}{dv}\right)^{u/2}\binom{{n\choose 2}}{dn/2}\left(\frac{3n}{vd}\right)^{v/2}\left(\frac{3}{\delta d}\right)^{(n-v)/2}}{g(\mathbf{d})}\right).
\end{gathered}\]


Using (8) we get that\par


\[\begin{gathered}
g(\mathbf{d})\geq\Omega(1)\binom{{n\choose 2}}{dn/2}((2+o(1))\pi d)^{-n/2}\geq\binom{{n\choose 2}}{dn/2}(7d)^{-n/2},
\end{gathered}\]


and so\par


\[\begin{gathered}
\displaystyle h(k,u,v,D) \displaystyle= \displaystyle O\left(n^{7/2}\ \left(\frac{4n}{dv}\right)^{u/2}\left(\frac{3n}{vd}\right)^{v/2}\left(\frac{3}{\delta d}\right)^{(n-v)/2}(7d)^{n/2}\right) \\
\displaystyle= \displaystyle O\left(n^{7/2}\ \left(\frac{4n}{dv}\right)^{u/2}\left(\frac{\delta n}{v}\right)^{v/2}\left(\frac{21}{\delta}\right)^{n/2}\right)\,.
\end{gathered}\]


Now, let $h(v):=(\delta n/v)^{v/2}$ and note that $h^{\prime}(v)=\frac{1}{2}(\delta n/v)^{v/2}(\log(\delta n/v)-1)$. Hence, $h(v)$ is maximized for $v=\delta n/e$ and we get that for any positive integer $v$,\par


\[\begin{gathered}
\left(\frac{\delta n}{v}\right)^{v/2}\leq\max_{v}h(v)=\exp\left(\frac{\delta n}{2e}\right)\leq 2^{n/2}\,,
\end{gathered}\]


since $\exp(\delta/e)\leq\exp(1/e)\approx 1.445\leq 2$. It follows that\par


\[\begin{gathered}
h(k,u,v,D)=O\left(n^{7/2}\ \left(\frac{4n}{dv}\right)^{u/2}\left(\frac{42}{\delta}\right)^{n/2}\right)=O\left(n^{7/2}\ \left(\left(\frac{4}{C}\right)^{\delta}\cdot\frac{42}{\delta}\right)^{n/2}\right)=O\left(5^{-n}\right),
\end{gathered}\]


provided that $C$ is large enough so that
$\left(\frac{4}{C}\right)^{\delta}\cdot\frac{42}{\delta}<1/5^{2}$.
Since $\delta$ is fixed, this condition can be easily satisfied and we may now finally define the constant $C$:\par


\[\begin{gathered}
C=C(\delta):=\max\left\{\frac{4}{(\delta/1200)^{1/\delta}},\frac{500}{\delta}\right\}.
\end{gathered}\]




We conclude that if $|V|\geq n^{3/4}$, then\par


\[\begin{gathered}
\mathbb{P}(\mathcal{E}_{k}(U,V))=e^{-\Omega(dn)}+e^{-\Omega(dn^{1/2})}+O\left(\sum_{D}h(k,u,v,D)\right)=O\left(n^{2}\,5^{-n}\right).
\end{gathered}\]


If $v<n^{3/4}$ then, as mentioned earlier, we do not get the term $\left(\frac{3n}{vd}\right)^{(v-1)/2}$ in the estimation of $h(k,u,v,D)$ (see (12)). However, for $v<n^{3/4}$, this term does not help us much anyway: $\left(\frac{3n}{vd}\right)^{(v-1)/2}=\exp(-\Theta(v\log v))\geq\exp(-n^{4/5})$. Hence, regardless of the size of $V$,\par


\[\begin{gathered}
\mathbb{P}(\mathcal{E}_{k}(U,V))=O\left(n^{2}\exp(n^{4/5})\,5^{-n}\right).
\end{gathered}\]


Finally, by the union bound,\par


\[\begin{gathered}
\mathbb{P}\left(\exists k,U,V\,\,\mathcal{E}_{k}(U,V)\right)\leq n\cdot 2^{n}\cdot 2^{n}\cdot O\left(n^{2}\exp(n^{4/5})\,5^{-n}\right)=o(1),
\end{gathered}\]


which finishes the proof of the theorem.
$\vartriangleleft$\par






\par\medskip\noindent\textbf{
4 Proof of Theorem 3: Sparse Case}\par




Sparser graphs clearly require more rounds of CR. Consider any $d$-regular graph with diameter $D$ and let $u$ and $v$ be any two vertices at distance $D$ from each other. CR run on the initial partition $V_{1}=\{v\}$ and $V_{2}=[n]\setminus\{v\}$ requires at least $D-2$ rounds to converge. Indeed, after $D-2$ rounds there are at least two vertices at distance at least $D-1$ from $v$ that are still of the same colour.\par



% SOURCE-ID Thmtheorem19


\par\medskip\noindent\textbf{
\textbf{Theorem 19}\textbf{.}
}\par



\emph{Let $d_{0}$ be large enough, let $d_{0}\leq d=o(n)$, and let $\mathbf{G}_{n}\sim\mathcal{G}_{n,d}$. Let $k\in\mathbb{N}$ be an arbitrary constant. Then, the following holds whp: for every non-trivial partition $[n]=V_{1}^{0}\sqcup V_{2}^{0}$ of the vertex set of $\mathbf{G}_{n}$, after at most $\operatorname{diam}(\mathbf{G}_{n})+2$ rounds of CR, there exists a partition $[n]=V_{1}\sqcup\ldots\sqcup V_{k}$ such that for any $i\in[k]$, $n/3k\leq|V_{i}|\leq 3n/k$ and $V_{i}$ is a union of some colour classes.}\par



% END-SOURCE-ID Thmtheorem19



\par\medskip\noindent\textbf{
4.1 Anti-concentration Results}\par



% SOURCE-ID Thmtheorem20


\par\medskip\noindent\textbf{
\textbf{Lemma 20}\textbf{.}
}\par



\emph{Let $\ell$ be a large enough constant, and let $d_{0}=d_{0}(\ell)$ be another large enough constant. Let $d_{0}\leq d=o(n)$ and $\mathbf{G}_{n}\sim\mathcal{G}_{n,d}$. Then, the following property holds whp: for every set $U$ of size $|U|\in[\frac{n\ell}{d},\frac{n}{2}]$ and every non-negative integer $s$, the number of vertices in $[n]\setminus U$ that have exactly $s$ neighbours in $U$ is at most $10n/\ln\ell$.}\par



% END-SOURCE-ID Thmtheorem20



\par\medskip\noindent\textbf{Proof.}\par




Due to Claim 10, the probability that there exists a set $U$ of size $m\in[n\ell/d,n/2]$ and a set $V\subset[n]\setminus U$ of size $t=10n/\ln\ell$ such that the number of edges between $U$ and $V$ is more than $dm/4$ is at most\par


\[\begin{gathered}
\displaystyle\sum_{m=n\ell/d}^{n/2}{n\choose m}{n\choose t}{mt\choose dm/4}\left(\frac{2d}{n}\right)^{dm/4} \displaystyle\leq\sum_{m=n\ell/d}^{n/2}\left(\frac{en}{t}\right)^{t}\left(\frac{en}{m}\cdot\left(\frac{4et}{d}\cdot\frac{2d}{n}\right)^{d/4}\right)^{m} \\
\displaystyle\leq(\ln\ell)^{t}\sum_{m=n\ell/d}^{n/2}\left(\frac{ed}{\ell}\cdot\left(\frac{300}{\ln\ell}\right)^{d/4}\right)^{m} \\
\displaystyle\leq n(\ln\ell)^{10n/\ln\ell}\left(1/2\right)^{(n\ell/d)d/5}=o(1).
\end{gathered}\]


Therefore, it suffices to prove the lemma for $s$ such that $10sn/\ln\ell<dm/4$. So, we may assume that $s<dm\ln\ell/(40n)$. On the other hand, by the Expander Mixing Lemma and Theorem 13, whp the number of edges between any set $U$ of size $m/2$ and any set $V\subset[n]\setminus U$ of size $t$ is at most $3td/4$. So, we may also assume that $s<3d/4$.\par




Next, by the Expander Mixing Lemma and Theorem 13, whp, for every set $U$ of size $m\in[n\ell/d,n/2]$ and every set $V\subset[n]\setminus U$ of size $t$, there are at least $md/3$ edges between $U$ and $[n]\setminus(U\cup V)$, and at most $2d(n-m-t)/3$ edges between $[n]\setminus(U\cup V)$ and $U\cup V$.\par




Fix a set $U$ of size $m\in[\frac{n\ell}{d},\frac{n}{2}]$ and a set $V\subset[n]\setminus U$ of size $t$. Fix a non-negative integer
$s\leq\min\left\{\frac{3d}{4},\frac{dm\ln\ell}{40n}\right\}$. Let us estimate the probability that every vertex from $V$ has exactly $s$ neighbours in $U$. Let us order the vertices in $V$ arbitrarily: $x_{1},x_{2},\ldots,x_{t}$, where $t=|V|=10n/\ln\ell$. Let $\mathcal{E}$ be the event that every $x_{i}$ has $s$ neighbours in $U$. Let $h_{in}\geq d(n-m-t)/6$ and $h_{out}\geq md/3$ be integers such that\par


\[\begin{gathered}
\mathbb{P}(\mathcal{E}\wedge\{|E(\mathbf{G}_{n}[[n]\setminus(V\cup U)]|=h_{in},\,|E(\mathbf{G}_{n}[U\times([n]\setminus(U\cup V))]|=h_{out}\})\text{ is maximum}.
\end{gathered}\]


Let $\Sigma_{0}$ be the set of all $d$-regular graphs $G$ on $[n]$ satisfying $\mathcal{E}$ and such that $G[[n]\setminus(V\cup U)]$ and $G[U\times([n]\setminus(U\cup V))]$ have exactly $h_{in}$ and $h_{out}$ edges, respectively. The following claim completes the proof of Lemma 20, see the proof in the extended version of the paper [29].\par



% SOURCE-ID Thmtheorem21


\par\medskip\noindent\textbf{
Claim 21. 
}\par



$\mathbb{P}(\mathbf{G}_{n}\in\Sigma_{0})\leq\ell^{-t/5}$.\par



% END-SOURCE-ID Thmtheorem21



Indeed, by the union bound over $U,V$ and the number of edges outside of $V$ ($h_{in},h_{out}$), we get that probability that there exist sets $V,U$ such that $\mathcal{E}$ holds is at most\par


\[\begin{gathered}
\displaystyle o(1)+(nd)^{2} \displaystyle\sum_{m=n\ell/d}^{n/2}{n\choose m}{n\choose t}\ell^{-t/5}\leq o(1)+n^{4}2^{n}\left(\frac{en}{t\ell^{1/5}}\right)^{t} \\
\displaystyle\leq o(1)+n^{4}2^{n}\left(\frac{\ln\ell}{\ell^{1/5}}\right)^{10n/\ln\ell}\leq o(1)+n^{4}2^{n}e^{-n}=o(1).
\end{gathered}\]


$\blacktriangleleft$\par



% SOURCE-ID Thmtheorem22


\par\medskip\noindent\textbf{
\textbf{Lemma 22}\textbf{.}
}\par



\emph{Let $\ell$ be large enough constant, and let $d_{0}=d_{0}(\ell)$ be another large enough constant. Let $d_{0}\leq d=o(n)$ and $\mathbf{G}_{n}\sim\mathcal{G}_{n,d}$. Then, the following property holds whp: for every set $U$ of size $|U|\in[\frac{n}{2d},\frac{n\ell}{d}]$ and every integer $s$ such that $1\leq s\leq\ell$, there are at most $0.999n$ vertices that have exactly $s$ neighbours in $U$.}\par



% END-SOURCE-ID Thmtheorem22




\par\medskip\noindent\textbf{Proof.}\par




By the Expander Mixing Lemma and Theorem 13, whp every set $V$ of size $n/4$ induces at most $dn/15$ edges. Let $\mathcal{E}$ be the event that there exists a set $V$ of size $n/4$ with more than $dn/15$ edges.\par




Let $\varepsilon=0.001$. Fix a set $U$ of size $m\in[\frac{n}{2d},\frac{\ell n}{d}]$ and a set $V=\{x_{1},\ldots,x_{t}\}\subset[n]\setminus U$ of size $n(1-\varepsilon)$. Divide $V=V^{\prime}\sqcup V^{\prime\prime}$, where $V^{\prime}$ consists of the first $n/4$ vertices. Expose edges inside $V^{\prime}$ and assume that $E:=E(V^{\prime})$ has size at most $dn/15$. Let $\tilde{V}^{\prime}\subset V^{\prime}$ be the set of vertices that have at most $d/2$ neighbours in $V^{\prime}$. Clearly, $|\tilde{V}^{\prime}|\geq n/12$. Without loss of generality, we assume $\tilde{V}^{\prime}=\{x_{1},\ldots,x_{t^{\prime}}\}$, where $t^{\prime}\geq n/12$.\par




Let $\Sigma_{0}$ be the set of $d$-regular graphs $G$ on $[n]$ such that $G[V^{\prime}]=E$ and each vertex in $V$ has exactly $s$ neighbours in $U$.
Let $\Sigma_{1}$ be the set of $d$-regular graphs $G$ on $[n]$ such that $G[V^{\prime}]=E$ and\par


\begin{itemize}


\item[] 
$\blacksquare$ 


$x_{1}$ has $s+1$ neighbours in $U$,\par





\item[] 
$\blacksquare$ 


each vertex $x_{2},\ldots,x_{t}$ has exactly $s$ neighbours in $U$, except for some $x_{i}\in V^{\prime\prime}$, whereas $x_{i}$ has $s-1$ neighbours in $U$,\par





\end{itemize}


We shall prove that $|\Sigma_{1}|\geq|\Sigma_{0}|/3$.
Take $G\in\Sigma_{0}$ and consider a tuple of vertices $(y,u,v)$ such that\par


\begin{itemize}


\item[] 
$\blacksquare$ 


$y\in V^{\prime\prime}$, $u\in U$, $v\in V^{\prime\prime}$,\par





\item[] 
$\blacksquare$ 


and $\{x_{1},y\},\{u,v\}\in E(G)$, $\{x_{1},u\},\{y,v\}\notin E(G)$.\par





\end{itemize}


If we switch\par


\[\begin{gathered}
\{x_{1},y\},\{u,v\}\mapsto\{x_{1},u\},\{y,v\},
\end{gathered}\]


we get a graph from $\Sigma_{1}$. For every $G\in\Sigma_{0}$ the number of \emph{forward switchings} is at least $(d/2-s)(((3/4-\varepsilon)n-d)s-sd)$. On the other hand, for every graph $G\in\Sigma_{1}$, the number of \emph{backward switchings} is at most $(s+1)(3n/4)d$. We get\par


\[\begin{gathered}
|\Sigma_{0}|(d/2-s)(((3/4-\varepsilon)n-d)s-sd)\leq|\Sigma_{1}|(s+1)(3n/4)d
\end{gathered}\]


implying $|\Sigma_{1}|\geq|\Sigma_{0}|/3$, as desired.\par




We now let $\Sigma_{2}$ be the set of $d$-regular graphs $G$ on $[n]$ such that\par


\begin{itemize}


\item[] 
$\blacksquare$ 


$|N_{U}(x_{1})|\in[s,s+1]$,\par





\item[] 
$\blacksquare$ 


$x_{2}$ has $s+1$ neighbours in $U$,\par





\item[] 
$\blacksquare$ 


each vertex $x_{3},\ldots,x_{t}$ has exactly $s$ neighbours in $U$, except for $x_{i_{1}}\in V^{\prime\prime}$ (when $|N_{U}(x_{1})|=s$) and $x_{i_{1}},x_{i_{2}}\in V^{\prime\prime}$ (when $|N_{U}(x_{1})|=s+1$) that have $s-1$ neighbours in $U$.\par





\end{itemize}


Take $G\in\Sigma_{0}\cup\Sigma_{1}$ and consider a tuple of vertices $(y,u,v)$ such that\par


\begin{itemize}


\item[] 
$\blacksquare$ 


$y\in V^{\prime\prime}$, $u\in U$, $v\in V^{\prime\prime}$, and $|N_{U}(v)|=s$,\par





\item[] 
$\blacksquare$ 


and $\{x_{2},y\},\{u,v\}\in E(G)$, $\{x_{2},u\},\{y,v\}\notin E(G)$.\par





\end{itemize}


If we switch as in (14), then
we get a graph from $\Sigma_{1}$. For every $G\in\Sigma_{0}\cup\Sigma_{1}$ the number of \emph{forward switchings} is at least $(d/2-s)(((3/4-\varepsilon)n-d-1)s-sd)$. On the other hand, for every $G\in\Sigma_{1}$, the number of \emph{backward switchings} is at most $(s+1)(3n/4)d$, as before. Thus\par


\[\begin{gathered}
|\Sigma_{0}\cup\Sigma_{1}|(d/2-s)(((3/4-\varepsilon)n-d-1)s-sd)\leq|\Sigma_{2}|(s+1)(3n/4)d
\end{gathered}\]


implying $|\Sigma_{2}|\geq|\Sigma_{0}\cup\Sigma_{1}|/3$, as well.\par




Similarly, we define $\Sigma_{3},\ldots,\Sigma_{n/12}$. For the $i$-th set $\Sigma_{i}$, we get that\par


\[\begin{gathered}
|\Sigma_{0}\cup\ldots\cup\Sigma_{i-1}|(d/2-s)(((3/4-\varepsilon)n-d-(i-1))s-sd)\leq|\Sigma_{i}|(s+1)(3n/4)d,
\end{gathered}\]


implying $|\Sigma_{i}|\geq|\Sigma_{0}\cup\ldots\cup\Sigma_{i-1}|/3$ for all $i$. In particular, we get\par


\[\begin{gathered}
\displaystyle|\Sigma_{7}| \displaystyle\geq\frac{|\Sigma_{0}|+|\Sigma_{1}|+\ldots+|\Sigma_{6}|}{3} \\
\displaystyle\geq\frac{|\Sigma_{0}|+|\Sigma_{0}|/3+(|\Sigma_{0}|+|\Sigma_{1}|)/3+\ldots+(|\Sigma_{0}|+\ldots+|\Sigma_{5}|)/3}{3} \\
\displaystyle>\frac{|\Sigma_{0}|+|\Sigma_{0}|/3+5(|\Sigma_{0}|+|\Sigma_{0}|/3)/3}{3}>\frac{7}{6}|\Sigma_{0}|.
\end{gathered}\]


In a similar way, for every $i$,
$|\Sigma_{7i+7}|\geq\frac{|\Sigma_{7i}|+\ldots+|\Sigma_{7i+6}|}{3}>\frac{7}{6}|\Sigma_{7i}|$.
Thus, we get $|\Sigma_{n/12}|>(7/6)^{n/84}|\Sigma_{0}|$,
implying $\mathbb{P}(\mathbf{G}_{n}\in\Sigma_{0}\mid E(\mathbf{G}_{n}[V^{\prime}])=E)<(7/6)^{-n/84}$. The union bound over $U$ and $V$ gives us that\par


\[\begin{gathered}
\mathbb{P}(\neg\mathcal{E})+{n\choose\varepsilon n}{n\choose\ell n/d}\sum_{E:\,|E|\leq dn/15}\mathbb{P}(\mathbf{G}_{n}\in\Sigma_{0}\mid E(\mathbf{G}_{n}[V^{\prime}])=E)\cdot\mathbb{P}(E(\mathbf{G}_{n}[V^{\prime}])=E)\\
=o(1)+e^{(\varepsilon\ln(e/\varepsilon)+o(1))n}(7/6)^{-n/84}=o(1),
\end{gathered}\]


which completes the proof of the lemma.
$\blacktriangleleft$\par






\par\medskip\noindent\textbf{
4.2 Proof of Theorem 19
}\par




With Lemmas 14, 20, and 22 at hand, we can easily prove Theorem 19.\par




\par\medskip\noindent\textbf{Proof of Theorem 19.}\par




Fix any $k\in\mathbb{N}$, and let $\ell=e^{30k}$. Let $c>0$ be a small enough constant as in Lemma 14. Let $d_{0}=d_{0}(\ell)$ be a large enough constant as in Lemmas 14, 20, and 22. Moreover, we will adjust constants $c$ or $d$, if needed, for some of the claims below to hold. Since we aim for the statement that holds whp, we may assume that the statements in Lemmas 14, 20, and 22 hold deterministically.\par




Consider any non-trivial partition $[n]=V_{1}^{0}\sqcup V_{2}^{0}$. Our goal is to show that after at most $\operatorname{diam}(\mathbf{G}_{n})+3$ many rounds of CR, we get a partition into colour classes that have sizes at most $n/3k$. To get the desired partition into $k$ parts, each of size at least $n/3k$ but at most $3n/k$, one can iteratively merge colour classes as we did in the proof of Theorem 15.\par




Let $u=\min\{|V_{1}^{0}|,|V_{2}^{0}|\}$ and let $U$ be a colour class of size $u$. Suppose first that $u<cn/d$. Let $r$ be the largest integer such that $ud(d-1)^{r-1}\leq cn$. We may adjust $c$ and $d$, if needed, to make sure that $1-100c-4\ln d/d\geq 1/2$, which, in particular, implies that $c\leq 1/200$. It follows from Lemma 14 that\par


\[\begin{gathered}
|S_{r}(U)|\geq(1-100c-4\ln d/d)ud(d-1)^{r-1}>\frac{cn}{2d},
\end{gathered}\]


and clearly $|S_{r}(U)|\leq ud(d-1)^{r-1}\leq cn\leq n/2$. After $r\leq\operatorname{diam}(\mathbf{G}_{n})-1$ rounds of CR, $S_{r}(U)$ is a union of some colour classes. We may merge them together at this point and continue the process from there.\par




Suppose now that $U$ is a colour class of size $u=|U|\in[\frac{cn}{2d},\frac{n}{2d}]$. Let $U^{\prime}$ be an arbitrary subset of $U$ of size $\frac{cn}{2d}$. On the one hand, trivially, $|S_{1}(U)|\leq d|U|\leq n/2$. On the other hand, it follows from Lemma 14 that $|S_{1}(U^{\prime})|\geq|U^{\prime}|d/2=cn/4$ which implies that $|S_{1}(U)|\geq|S_{1}(U^{\prime})|-|U|\geq cn/4-n/2d\geq n\ell/d$, when $d$ is large enough. Since $S_{1}(U)$ is a union of some colour classes, we may merge them into one large class and continue from there.\par




Suppose this time that $U$ is a colour class of size $u=|U|\in[\frac{n}{2d},\frac{n\ell}{d}]$. We may adjust $d$, if needed, to make sure $u\leq\frac{n\ell}{d}\leq\frac{n}{2(\ell+1)}$. After one round of CR, $[n]\setminus U$ is partitioned into sets $W_{i}$ ($i\in\mathbb{N}\cup\{0\}$); set $W_{i}$ consists of vertices with exactly $i$ neighbours in $U$. Let $A=\bigcup_{i\leq\ell}W_{i}$ and let $B=\bigcup_{i\geq\ell+1}W_{i}$. Clearly, $|U|+|A|+|B|=n$. Note that, on the one hand, the number of edges between $U$ and its complement is at least $|B|(\ell+1)$. On the other hand, it is trivially at most $|U|d\leq n\ell$. We conclude that
$|B|\leq\frac{\ell}{\ell+1}n=\left(1-\frac{1}{\ell+1}\right)n$,
and so\par


\[\begin{gathered}
|A|=n-|B|-|U|\geq\frac{n}{\ell+1}-u\geq\frac{n}{2(\ell+1)}.
\end{gathered}\]




Our goal is to show that one can always merge some sets $W_{i}$ together to get a colour class of size at most $n/2$ but at least $\frac{n}{2(\ell+1)}$, which is at least $\frac{n\ell}{d}$, provided that $d$ is large enough. To that end, we will consider a few cases. If $|A|\leq n/2$, then we can simply take the entire set $A$ for the desired colour class. If $|A|>n/2$ but $|A|\leq(1-1/(\ell+1))n$, then we may take the entire set $B$ since $|B|=n-|A|-|U|\geq n/(\ell+1)-u\geq n/2(\ell+1)$ and, trivially, $|B|=n-|A|-|U|<n/2$.\par




It remains to concentrate on the case when $|A|\geq(1-1/(\ell+1))n$. Suppose first that $|W_{i}|\geq n/2$ for some $0\leq i\leq\ell$. It follows from Lemma 22 that $|W_{i}|\leq 0.999n$. Then, we can take $A\setminus W_{i}$ for the desired colour class since, trivially, $|A\setminus W_{i}|\leq n-|W_{i}|\leq n/2$ and\par


\[\begin{gathered}
|A\setminus W_{i}|\geq|A|-|W_{i}|\geq 0.001n-\frac{n}{\ell+1}\geq\frac{n\ell}{d},
\end{gathered}\]


provided that $d$ is large enough. If $n/4\leq|W_{i}|<n/2$ for some $0\leq i\leq\ell$, then we may simply take $W_{i}$ as our colour class. Suppose then that $|W_{i}|<n/4$ for all $0\leq i\leq\ell$. Then, we may start with set $A$ and remove $W_{i}$’s, one by one, and at some point we get a set of size at most $n/2$ but at least $n/4$.\par




Finally, suppose that $U$ is a colour class of size $u=|U|\in[\frac{n\ell}{d},\frac{n}{2}]$. It follows immediately from Lemma 20 that after one round of CR, the complement of $U$ is partitioned into sets of size at most $10n/\ln\ell=n/3k$. We can group some of them together to get a colour class of size at least $n/3k\geq n\ell/d$ but at most $2n/3k\leq n/2$ to make sure that after one more round $U$ is also partitioned into sets of size at most $n/3k$. This completes the proof.
$\blacktriangleleft$\par







\par\medskip\noindent\textbf{
5 Proof of Theorem 4: Dense Case}\par




Here, we prove the following.\par



% SOURCE-ID Thmtheorem23


\par\medskip\noindent\textbf{
\textbf{Theorem 23}\textbf{.}
}\par



\emph{Let $\varepsilon>0$, let $n^{1/2+\varepsilon}\leq d\leq n/2$, and let $\mathbf{G}_{n}\sim\mathcal{G}_{n,d}$.
There exists some universal large constant $k\in\mathbb{N}$ such that the following holds whp: for every partition $[n]=V_{1}\sqcup\ldots\sqcup V_{k}$ of the vertex set of $\mathbf{G}_{n}$ such that for any $i\in[k]$, $n/3k\leq|V_{i}|\leq 3n/k$, after three rounds of CR, there are only singleton colour classes.}\par



% END-SOURCE-ID Thmtheorem23



We start from a simple auxiliary lemma, the proof of which is in the extended version of the paper [29].\par



% SOURCE-ID Thmtheorem24


\par\medskip\noindent\textbf{
\textbf{Lemma 24}\textbf{.}
}\par



\emph{Let $\varepsilon>0$, $n^{1/2+\varepsilon}\leq d=d(n)\leq n/2$, and let $\mathbf{G}_{n}\sim\mathcal{G}_{n,d}$.
For every pair of vertices $u,v\in[n]$, let $M_{u,v}\subset N(u)$ and $M^{\prime}_{u,v}\subset N(v)\setminus N(u)$ be sets of size $\lfloor n/(20d)\rfloor$ chosen uniformly at random. Then the following events hold whp for any pair of vertices $u,v$ in $\mathbf{G}_{n}$:}\par


\begin{enumerate}


\item[] 
1. 


$|N(u)\cap N(v)|<2d/3$\emph{;}\par





\item[] 
2. 


$|N(u)\cup N(v)\cup\{u,v\}|<4n/5$\emph{;}\par





\item[] 
3. 


$|N(M^{\prime}_{u,v})\setminus(N(\{u,v\})\cup N(M_{u,v}))|>n/25$\emph{ whenever }$d\leq n/20$\emph{.}\par





\end{enumerate}



% END-SOURCE-ID Thmtheorem24



\par\medskip\noindent\textbf{Proof of Theorem 23.}\par




Due to the Expander Mixing Lemma and Theorem 13, whp between any set $N$ of size $\Omega(d)$ and any set $W$ of size $\Omega(n)$, there are $(1\pm o(1))|N||W|\frac{d}{n}$ edges. We denote the intersection of this event with the event from the assertion of Lemma 24 by $\mathcal{E}$.\par




Suppose that $k$ is as large as needed, and fix any partition $[n]=V_{1}\sqcup\ldots\sqcup V_{k}$ such that each part has size in the range $[n/3k,3n/k]$ as in the statement of the theorem. For each $i\in[k]$, define $d_{i}:[n]\to\mathbb{Z}$ so that $d_{i}(w)=|V_{i}\cap N(w)|$. Finally, for each vertex $u\in[n]$, we interpret $c_{i}(u)$ as the colour of $u$ after $i$ rounds of CR. For our goal, it suffices to show that whp
no two vertices have the same value of $c_{3}(\cdot)$.\par




We proceed as follows. Fix a pair of vertices $u,v\in[n]$ and expose the neighbourhoods of $u$ and $v$.
Note that $c_{3}(u)=c_{3}(v)$ if and only if there exists a bijection $b:N(v)\mapsto N(u)$ such that for any $w\in N(v)$ we have $c_{2}(b(w))=c_{2}(w)$. Fix such a bijection (in $d!$ ways). Define $N^{\prime}:=N(v)\setminus(N(u)\cup\{u\})$.
Choose arbitrarily a set $M^{\prime}\subset N^{\prime}$ of $\lfloor n/(20d)\rfloor$ vertices from $N^{\prime}$, and let $M=b(M^{\prime})\subset N(u)$. Expose all edges that touch $M\cup N^{\prime}$. Let\par


\[\begin{gathered}
M^{\prime\prime}=N^{\prime}\cup N(M^{\prime})\setminus(N(u)\cup N(M)\cup\{u,v\}).
\end{gathered}\]


Due to symmetry we may assume $|N(v)\cup N(M)|\geq|N(u)\cup N(M^{\prime})|$ and we extend the bijection $b$ to an injection $b:N(v)\cup N(M)\mapsto N(u)\cup N(M^{\prime})$ such that, for every $w\in M$ and every $w^{\prime}\in N(w)$, we get $c_{1}(b(w^{\prime}))=c_{1}(w^{\prime})$ and $b(w^{\prime})\in N(b(w))$. The number of ways to define such an extension is at most $(d!)^{n/20d}$.\par




Let $W:=[n]\setminus(\{u,v\}\cup N(u)\cup N(v)\cup N(M)\cup N(M^{\prime}))$, and $W_{i}:=W\cap V_{i}$. After exposing every edge except those between $M^{\prime\prime}$ and $W$, we can determine the values of $d_{i}(w)$ for every $w\in N(M)$ (as every neighbour of every vertex in $N(M)$ has been exposed). Therefore, the injection $b$ also identifies $d_{i}(w)$ for every $i\in[k]$ and every $w\in M^{\prime\prime}$. Let $S$ denote the number of pairs $(i,w)$ with $1\leq i\leq k$ and $w\in M^{\prime\prime}$ such that $|W_{i}|\geq n/(30k)$ and $|W_{i}|d/(2n)\leq|N(w)\cap W_{i}|\leq 3|W_{i}|d/(2n)$. The number of ways to choose the remaining neighbours of the vertices in $M^{\prime\prime}$ is\par


\[\begin{gathered}
\prod_{i=1}^{k}\prod_{w\in M^{\prime\prime}}\binom{|W_{i}|}{d_{i}(w)}\leq\left(\frac{960k}{d}\right)^{S/2}\binom{|W||M^{\prime\prime}|}{dn/2-m},
\end{gathered}\]


where $m$ is the number of exposed edges.
Indeed for any positive integers $a_{1},a_{2},b_{1},b_{2}$ with $b_{1}\geq b_{2}$ and $b_{i}\leq 3a_{i}/4$ for $i=1,2$ we have by Corollary 6 that\par


\[\begin{gathered}
\displaystyle\binom{a_{1}}{b_{1}}\binom{a_{2}}{b_{2}} \displaystyle\leq\frac{2}{3}\sqrt{\frac{(b_{1}+b_{2})(a_{1}+a_{2})a_{1}a_{2}}{(a_{1}+a_{2})b_{1}(a_{1}-b_{1})b_{2}(a_{2}-b_{2})}}\binom{a_{1}+a_{2}}{b_{1}+b_{2}} \\
\displaystyle\leq\frac{8}{3}\sqrt{\frac{b_{1}+b_{2}}{b_{1}b_{2}}}\binom{a_{1}+a_{2}}{b_{1}+b_{2}}\leq{\frac{4}{\sqrt{b_{2}}}}\binom{a_{1}+a_{2}}{b_{1}+b_{2}}.
\end{gathered}\]




Let us show that the event $\mathcal{E}$ implies $S\geq kn/1100$. Indeed, this event implies that $|W|\geq n/5$ (if $d>n/20$, then $W=[n]\setminus(N(\{u,v\})\cup\{u,v\})$ and has size at least $n/5$ by the second assertion of Lemma 24; if $d\leq n/20$, then $|W|\geq n-2d-2(n/(20d))d>n/5$). Therefore, there are at least $n/6$ vertices in the union of $W_{i}$ such that $|W_{i}|\geq n/(30k)$. Thus, there are at least $(n/6)/(3n/k)=k/18$ such $W_{i}$. Fix such a $W_{i}$. Since $\mathcal{E}$ holds, any subset $\tilde{N}\subset M^{\prime\prime}$ of size $\Omega(n)$ sends $(1\pm o(1))|\tilde{N}||W_{i}|$ edges to $W_{i}$. Moreover, $|M^{\prime\prime}|\geq n/60$. Indeed, if $d>n/20$, then $M^{\prime\prime}=N^{\prime}$ which has size at least $d/3>n/60$ by the first assertion of Lemma 24; if $d\leq n/20$, then $|M^{\prime\prime}|>n/25$ by the third assertion. The event $\mathcal{E}$ also implies that number of vertices in $M^{\prime\prime}$ that have less than $|W_{i}|\frac{d}{2n}$ or more than $|W_{i}|\frac{3d}{2n}$ edges in $W_{i}$, is $o(n)$. So, indeed $S>(1-o(1))(k/18)(n/60)>kn/1100$.\par




Using (8) and letting $g(\mathbf{d})=(d,\ldots,d)\in\mathbb{Z}^{n}$, the probability that there exists $u,v\in[n]$ and a partition $V_{1}\sqcup\ldots\sqcup V_{k}$ such that $c(u)=c(v)$ is then at most\par


\[\begin{gathered}
\mathbb{P}(\neg\mathcal{E})+\frac{1}{g(\mathbf{d})}n^{2}2^{kn}d!\sum_{N^{\prime}\subseteq[n]}\sum_{W\subseteq[n]\setminus N^{\prime}}\sum_{m=0}^{dn/2}\binom{\binom{n}{2}-|N^{\prime}||W|}{m}\left(\frac{960k}{d}\right)^{S/2}\binom{|W||N^{\prime}|}{dn/2-m}\\
=o(1)+d^{\Theta(n)}\left(\frac{960k}{d}\right)^{kn/2200}=o(1),
\end{gathered}\]


when $k$ is sufficiently large. This completes the proof of Theorem 23.
$\blacktriangleleft$\par






\par\medskip\noindent\textbf{
6 Proof of Theorem 4: Sparse Case}\par




Here we prove the following.\par



% SOURCE-ID Thmtheorem25


\par\medskip\noindent\textbf{
\textbf{Theorem 25}\textbf{.}
}\par



\emph{There exists a universal constant $k$ such that the following holds. Let $d_{0}=d_{0}(k)$ be large enough, let $d_{0}\leq d\leq n^{10/17}$, and let $\mathbf{G}_{n}\sim\mathcal{G}_{n,d}$.
Then whp: for every partition $[n]=V_{1}\sqcup\ldots\sqcup V_{k}$ of the vertex set of $\mathbf{G}_{n}$ such that for any $i\in[k]$, $n/3k\leq|V_{i}|\leq 3n/k$, after $\operatorname{diam}(\mathbf{G}_{n})+1$ rounds of CR, there are only singleton colour classes.}\par



% END-SOURCE-ID Thmtheorem25



We will use two direct corollaries of Lemma 14 from Section 4 in this proof. We first state these in Section 6.1, and then prove Theorem 25 in Section 6.2.\par




\par\medskip\noindent\textbf{
6.1 Sizes of Balls}\par




As in Lemma 14, let $c>0$ be small enough and $d_{0}$ be large enough. Let $d_{0}\leq d=o(n)$ and let $\mathbf{G}_{n}\sim\mathcal{G}_{n,d}$. The following two lemmas are direct corollaries of Lemma 14.\par



% SOURCE-ID Thmtheorem26


\par\medskip\noindent\textbf{
\textbf{Lemma 26}\textbf{.}
}\par



\emph{Whp, for every $r$ such that $d(d-1)^{r-1}\leq cn$, the following holds}\par


\begin{itemize}


\item[] 
$\blacksquare$ 


\emph{for every vertex }$u$\emph{,}\par


\[\begin{gathered}
|S_{r}(u)|\geq\left(1-100c-\frac{4\ln d}{d}\right)d(d-1)^{r-1};
\end{gathered}\]





\item[] 
$\blacksquare$ 


\emph{for every pair of vertices }$u\neq v$\emph{,}\par


\[\begin{gathered}
|S_{r}(v)\setminus B_{r}(u)|\geq|S_{r}(\{u,v\})|-|B_{r}(u)|\geq\left(1-200c-\frac{8\ln d}{d}-\frac{1}{d-2}\right)d(d-1)^{r-1}.
\end{gathered}\]





\end{itemize}



% END-SOURCE-ID Thmtheorem26



\par\medskip\noindent\textbf{Proof.}\par




The first assertion is just Lemma 14 applied with $U=\{u\}$. The second follows from $|U|=2$ in Lemma 14 together with the basic bound $|B_{r}(u)|\leq\sum_{i\leq r-1}d(d-1)^{i}$.
$\blacktriangleleft$\par



% SOURCE-ID Thmtheorem27


\par\medskip\noindent\textbf{
\textbf{Lemma 27}\textbf{.}
}\par



\emph{Whp}\par


\begin{itemize}


\item[] 
$\blacksquare$ 


\emph{for every set }$U$\emph{ of size }$\frac{cn}{d}$\emph{, there are at least }$(1-4\ln d/d-100c)cn$\emph{ vertices that have a neighbour in }$U$\emph{;}\par





\item[] 
$\blacksquare$ 


\emph{for any two disjoint sets }$U,V$\emph{ of size }$cn/d$\emph{, the number of vertices that have neighbours both in }$U$\emph{ and in }$V$\emph{ is at most }$|N(U)|/10$\emph{.}\par





\end{itemize}



% END-SOURCE-ID Thmtheorem27




\par\medskip\noindent\textbf{Proof.}\par




The first assertion follows immediately from Lemma 14 applied with $r=1$. The second assertion follows as well since whp for any two disjoint sets $U$ and $V$, the number of vertices that have neighbours in both sets is at most\par


\[\begin{gathered}
\displaystyle|N(U)|+|N(V)| \displaystyle-|N(U\cup V)|\leq 2d(cn/d)-(1-4\ln d/d-100c)2cn \\
\displaystyle=(4\ln d/d+100c)2cn\leq\frac{1}{10}(1-4\ln d/d-100c)cn\leq|N(U)|/10.\
\end{gathered}\]


$\blacktriangleleft$\par






\par\medskip\noindent\textbf{
6.2 Colour Refinement Run on a Vertex-coloured Random Graph}\par




Let $d$ be large enough. In what follows we assume that properties from Lemma 26 and Lemma 27 hold in $\mathbf{G}_{n}$ deterministically.\par




Fix a partition $[n]=V_{1}\sqcup\ldots\sqcup V_{k}$ as in the statement of the theorem. Assign to every vertex $x$ the colour $C_{0}(x)$ that equals the index of the set $V_{i}$ to which $x$ belongs. Let $\mathbf{D}$ be the diameter of $\mathbf{G}_{n}$. Consider the output $C_{t}$ of $t:=\mathbf{D}+1$ rounds of CR at the coloured graph. We want to prove that $C(u)\neq C(v)$ for any two different vertices $u,v\in[n]$.\par




Fix two vertices $u\neq v$. Assume $C_{t}(u)=C_{t}(v)$. Then, for every neighbour $x$ of $v$, there exists a neighbour $y$ of $u$ such that $C_{t-1}(x)=C_{t-1}(y)$. More generally, we have the following.\par



% SOURCE-ID Thmtheorem28


\par\medskip\noindent\textbf{
Claim 28. 
}\par



Let $r\in[t]$. For every vertex $a$ and every vertex $b$ such that $C_{t-r}(a)=C_{t-r}(b)$, and every neighbour $x$ of $a$, there exists a neighbour $y$ of $b$ such that $C_{t-r-1}(x)=C_{t-r-1}(y)$.\par



% END-SOURCE-ID Thmtheorem28



Let $r=\lfloor\log_{d-1}(\varepsilon n)\rfloor$, where $\varepsilon>0$ is a small enough constant. By Lemma 26, we have that\par


\[\begin{gathered}
\displaystyle|S_{r}(u)\setminus B_{r}(v)| \displaystyle\geq\left(1-200c-\frac{8\ln d}{d}-\frac{1}{d-2}\right)d(d-1)^{r-1} \\
\displaystyle>\frac{d}{(d-1)^{2}}\varepsilon n\left(1-200c-\frac{8\ln d}{d}-\frac{1}{d-2}\right)>\frac{\varepsilon}{2d}\cdot n.
\end{gathered}\]




Due to Claim 28, for every vertex $x\in S_{1}(u)\setminus B_{1}(v)$, there exists a vertex $f(x)\in B_{1}(v)$ such that $C_{t-1}(x)=C_{t-1}(f(x))$. Next, for every vertex $x\in S_{2}(u)\setminus B_{2}(v)$, let $\pi(x)\in S_{1}(u)\setminus B_{1}(v)$ be one of its “parents”. We have that $C_{t-1}(\pi(x))=C_{t-1}(f(\pi(x))$. Therefore, by Claim 28, there exists $f(x)\in N(f(\pi(x)))\subset B_{2}(v)$ such that $C_{t-2}(x)=C_{t-2}(f(x))$. We then define $f:B_{r}(u)\setminus B_{r}(v)\to B_{r}(v)$ by induction: for every $2\leq i\leq r$, assuming that $f$ has been defined on $B_{i-1}(u)\setminus B_{i-1}(v)$, and for every $x\in S_{i}(u)\setminus B_{i}(v)$, find its “parent” $\pi(x)\in S_{i-1}(u)\setminus B_{i-1}(v)$ and take $f(x)\in N(f(\pi(x))$ such that $C_{t-i}(x)=C_{t-i}(f(x))$.\par




Take $U\subset S_{r}(u)\setminus B_{r}(v)$ of size $\frac{\varepsilon n}{2d}$ and let $U^{\prime}:=f(U)\subset B_{r}(v)$. We have $|U^{\prime}|\leq|U|$. Without loss of generality, we assume $|U^{\prime}|=|U|$ (otherwise, we can extend $U^{\prime}$ arbitrarily to keep the two sets disjoint, and the argument below will still work). Note that $B_{r}(u)=B_{r-1}(u)\cup N(S_{r-1}(u))$, that
$|S_{r-1}(u)|\leq d(d-1)^{r-2}\leq\frac{d}{(d-1)^{2}}\varepsilon n<1.1\varepsilon\frac{n}{d},$
and that $U$ and $S_{r-1}(u)$ are disjoint. The same facts hold for $B_{r}(v)$. In particular,
$|S_{r-1}(u)\cup S_{r-1}(v)|<3\varepsilon\frac{n}{d}$.
Therefore, by the conclusion of Lemma 27, we have that\par


\[\begin{gathered}
\displaystyle|N(U)\setminus(N(U^{\prime})\cup B_{r}(u)\cup B_{r}(v))| \displaystyle>|N(U)|-\frac{7}{10}|N(U)|-|B_{r-1}(u)|-|B_{r-1}(v)| \\
\displaystyle>\frac{3}{10}(1-4\ln d/d-100c)\frac{1}{2}\varepsilon n-2\frac{d}{d-2}(d-1)^{r-1} \\
\displaystyle>0.14\cdot\varepsilon n-2\frac{d\varepsilon n}{(d-2)(d-1)}>0.1\cdot\varepsilon n.
\end{gathered}\]


Let $\mathcal{N}$ be a subset of $N(U)\setminus(N(U^{\prime})\cup B_{r}(u)\cup B_{r}(v))$ of size $\varepsilon n/10$.\par




We then extend $f$ to $\mathcal{N}$: Each vertex $x\in\mathcal{N}$ has $f(x)\in N(U^{\prime})$. Note that the set $X:=[n]\setminus(B_{r}(u)\cup B_{r}(v)\cup\mathcal{N}\cup f(\mathcal{N}))$
has size at least\par


\[\begin{gathered}
n-2\frac{d}{d-2}(d-1)^{r}-\frac{\varepsilon n}{2d}\cdot(2d)>n-2\frac{d}{d-2}\varepsilon\cdot n-\varepsilon\cdot n\geq n(1-4\varepsilon).
\end{gathered}\]


The set $X$ is partitioned into $k$ sets $X=V^{\prime}_{1}\sqcup\ldots\sqcup V^{\prime}_{k}$ so that $n(\frac{1}{3k}-4\varepsilon)\leq|V^{\prime}_{i}|\leq n\cdot\frac{3}{k}$.\par




Due to Claim 10, whp any set of size at most $3\varepsilon n$ induces at most $\frac{1}{100}\varepsilon dn$ edges:\par


\[\begin{gathered}
{n\choose 3\varepsilon n}{9\varepsilon^{2}n^{2}/2\choose\frac{1}{100}\varepsilon dn}\left(\frac{2d}{n}\right)^{\frac{1}{100}\varepsilon dn}\leq\left(\left(\frac{e}{3\varepsilon}\right)^{300}\left(900e\varepsilon\right)^{d}\right)^{\frac{1}{100}\varepsilon n}=o(1),
\end{gathered}\]


since $d$ is large and $\varepsilon$ is small enough. In particular, we may assume that there are at most $\frac{1}{100}\varepsilon dn$ edges between $\mathcal{N}$ and $\mathcal{N}\cup f(\mathcal{N})\cup B_{r}(u)\cup B_{r}(v)$. We get that there exists a subset $\mathcal{N}_{0}\subset\mathcal{N}$ of size $\frac{\varepsilon n}{50}$ such that each vertex in this set sends at least $\frac{3}{4}d$ edges to $X$.\par




Note that, for any vertex $x\in\mathcal{N}_{0}$, the equality $C_{t-r-1}(x)=C_{t-r-1}(f(x))$ implies $C_{1}(x)=C_{1}(f(x))$. Therefore, as soon as the sets $B_{r}(u),B_{r}(v)$ are exposed, the set $U$ is chosen, the sets $N(U),N(U^{\prime}),N(N(U^{\prime}))$ are exposed, and the set $\mathcal{N}_{0}$ is chosen, there should exist a function $f$ defined as above, that identifies the values of $|N_{V^{\prime}_{j}}(x)|$ for every $x\in\mathcal{N}_{0}$ and $j\in[k]$.\par




Therefore, we run the following exploration process of the random graph. First, we expose $B_{r}(u),B_{r}(v)$ and then choose $U\subset S_{r}(u)\setminus B_{r}(v)$ of size $\frac{\varepsilon n}{2d}$ arbitrarily. We then expose $N(U)$, $N(U^{\prime})$, and $N(N(U^{\prime}))$. We choose $f$ on $(B_{r}(u)\setminus B_{r}(v))\cup\mathcal{N}$ in at most $d^{2\varepsilon n}$ ways, since\par


\[\begin{gathered}
|(B_{r}(u)\setminus B_{r}(v))\cup\mathcal{N}|\leq|B_{r}(u)|+|N(U)|\leq\frac{d}{d-2}(d-1)^{r}+\frac{\varepsilon n}{2}\leq\frac{d}{d-2}\varepsilon n+\frac{\varepsilon n}{2}<2\varepsilon n.
\end{gathered}\]


Finally, we choose any set $\mathcal{N}_{0}\subset\mathcal{N}$ of size $\frac{\varepsilon n}{50}$ such that each vertex in this set sends at least $\frac{3}{4}d$ edges to $X$.\par




By the Expander Mixing Lemma and Theorem 13, whp between any two disjoint sets of size at least $n/(4k)$ and $n/2$, there are at least $dn/(10k)$ edges, and every set of size at least $n/2$ induces at least $dn/10$ edges.\par




Recall that every $x\in\mathcal{N}_{0}$ has a prescribed number of neighbours $g_{j}(x)$ in the set $V^{\prime}_{j}$.
By the Expander Mixing Lemma and Theorem 13, whp the number of edges between any two disjoint sets $U$ and $V$ of sizes $\Theta(n)$ equals $|U||V|d(1\pm\varepsilon)/n$. Therefore, for every set $V^{\prime}_{j}$, there exists a subset $\mathcal{N}^{\prime}_{j}\subset\mathcal{N}_{0}$ of size $\varepsilon n/100$ such that every $x\in\mathcal{N}^{\prime}_{j}$ has $g_{j}(x)\in[d/(10k),10d/k]$.\par




Let us estimate the probability that for every $j\in[k]$, every vertex from $\mathcal{N}^{\prime}_{j}$ has $g_{j}(x)$ neighbours in $V^{\prime}_{j}$. For every $j$, we order arbitrarily the vertices in $\mathcal{N}^{\prime}_{j}$: $x^{j}_{1},\ldots,x^{j}_{t}$, where
$t=\varepsilon n/100$. Let $\mathcal{E}$ be the event that, for every $j\in[k]$, every $x^{j}_{i}$ has $g_{j}(x^{j}_{i})$ neighbours in $V^{\prime}_{j}$. Let $h^{j}_{in}\geq dn/10$ and $h^{j}_{out}\geq dn/(10k)$ be integers such that\par


\[\begin{gathered}
\mathbb{P}\left(\mathcal{E}\wedge\bigwedge_{j=1}^{k}\left\{|E(\mathbf{G}_{n}[X\setminus V^{\prime}_{j}]|=h^{j}_{in},\,|E(\mathbf{G}_{n}[V^{\prime}_{j}\times(X\setminus V^{\prime}_{j})]|=h^{j}_{out}\right\}\right)\text{ is maximum}.
\end{gathered}\]


Let $\Sigma_{0}$ be the set of all $d$-regular graphs $G$ on $[n]$ satisfying $\mathcal{E}$ and such that, for all $j\in[k]$, $G[X\setminus V^{\prime}_{j}]$ and $G[V^{\prime}_{j}\times(X\setminus V^{\prime}_{j})]$ have exactly $h^{j}_{in}$ and $h^{j}_{out}$ edges, respectively\footnote{Sets $X$ and $V^{\prime}_{j}$ depend on $G$: given a graph $G$, we expose the balls around $u$ and $v$, which identify these sets. In what follows, we will perform switching operations on $G$ that preserve the exposed balls and, therefore, sets $X$ and $V^{\prime}_{j}$.}. The following claim completes the proof of Lemma 20 (its proof appears in the extended version of the paper [29]).\par



% SOURCE-ID Thmtheorem29


\par\medskip\noindent\textbf{
Claim 29. 
}\par



$\mathbb{P}(\mathbf{G}_{n}\in\Sigma_{0})\leq(k/d)^{\varepsilon nk/500}$.\par



% END-SOURCE-ID Thmtheorem29



Indeed, Claim 29 implies that $\mathbb{P}(\mathcal{E})\leq(dn)^{2k}(k/d)^{\varepsilon nk/500}$. Therefore, by the union bound\par


\[\begin{gathered}
\mathbb{P}(C_{t}(u)=C_{t}(v))\leq d^{2\varepsilon n}(dn)^{2k}(k/d)^{\varepsilon nk/500}=e^{-\Omega(kn)}
\end{gathered}\]


when $k$ is large enough and $d\gg k$. The union bound over the choice of partition $V_{1}\sqcup\ldots\sqcup V_{k}$ and over all pairs of distinct vertices $u,v$ completes the proof of Theorem 25.\par






\par\medskip\noindent\textbf{References}\par



\begin{itemize}


\item[] 
[1]

J. Alman, R. Duan, V. Vassilevska Williams, Y. Xu, Z. Xu, and R. Zhou.


More asymmetry yields faster matrix multiplication.


In \emph{Proceedings of the 2025 Annual ACM-SIAM Symposium on Discrete
Algorithms (SODA’25)}, pages 2005–2039, 2025.





\item[] 
[2]

N. Alon.


Eigenvalues and expanders.


\emph{Combinatorica}, 6(2):83–96, 1986.


doi:10.1007/BF02579166.





\item[] 
[3]

N. Alon and F. R. K. Chung.


Explicit construction of linear sized tolerant networks.


\emph{Discrete Mathematics}, 72:15–19, 1988.


doi:10.1016/0012-365X(88)90189-6.





\item[] 
[4]

N. Alon, M. Krivelevich, and V. H. Vu.


On the concentration of eigenvalues of random symmetric matrices.


\emph{Israel Journal of Mathematics}, 131:259–267, 2002.





\item[] 
[5]

M. Anastos, M. Kwan, and B. Moore.


Smoothed analysis for graph isomorphism.


In \emph{Proceedings of the 57th Annual ACM Symposium on Theory of
Computing (STOC’25)}, pages 2098–2106, 2025.





\item[] 
[6]

L. Babai.


Graph isomorphism in quasipolynomial time.


In \emph{Proceedings of the 48th Annual ACM Symposium on Theory of
Computing (STOC’16)}, pages 684–697, 2016.





\item[] 
[7]

L. Babai.


Canonical form for graphs in quasipolynomial time: preliminary
report.


In \emph{Proceedings of the 51st Annual ACM SIGACT Symposium on Theory
of Computing (STOC’19)}, pages 1237–1246, 2019.





\item[] 
[8]

L. Babai, P. Erdős, and S. M. Selkow.


Random graph isomorphism.


\emph{SIAM Journal on Computing}, 9(3):628–635, 1980.


doi:10.1137/0209047.





\item[] 
[9]

L. Babai and E. M. Luks.


Canonical labeling of graphs.


In \emph{Proceedings of the 15th Annual ACM Symposium on Theory of
Computing (STOC’83)}, pages 171–183, 1983.





\item[] 
[10]

R. Bauerschmidt, J. Huang, A. Knowles, and H.-T. Yau.


Edge rigidity and universality of random regular graphs of
intermediate degree.


\emph{Geometric and Functional Analysis}, 30:693–769, 2020.





\item[] 
[11]

N. Behague, D. Il’ković, and R. Montgomery.


A proof of the Kim-Vu sandwich conjecture.


\emph{arXiv preprint}, 2025.


arXiv:2510.20765.





\item[] 
[12]

C. Berkholz, P. Bonsma, and M. Grohe.


Tight lower and upper bounds for the complexity of canonical colour
refinement.


\emph{Theory of Computing Systems}, 60:581–614, 2017.


doi:10.1007/S00224-016-9686-0.





\item[] 
[13]

B. Bollobás.


Distinguishing vertices of random graphs.


\emph{Annals of Discrete Mathematics}, 13:33–50, 1982.





\item[] 
[14]

B. Bollobás.


\emph{Random Graphs}.


Cambridge University Press, 2 edition, 2001.





\item[] 
[15]

A. Z. Broder, A. M. Frieze, S. Suen, and E. Upfal.


Optimal construction of edge-disjoint paths in random graphs.


\emph{SIAM Journal on Computing}, 28(2):541–573, 1999.


doi:10.1137/S0097539795290805.





\item[] 
[16]

N. Chiba and T. Nishizeki.


Arboricity and subgraph listing algorithms.


\emph{SIAM Journal on Computing}, 14(1):210–223, 1985.


doi:10.1137/0214017.





\item[] 
[17]

N. A. Cook, L. Goldstein, and T. Johnson.


Size biased couplings and the spectral gap for random regular graphs.


\emph{Annals of Probability}, 46(1):72–125, 2018.





\item[] 
[18]

T. Czajka and G. Pandurangan.


Improved random graph isomorphism.


\emph{Journal of Discrete Algorithms}, 6:85–92, 2008.


doi:10.1016/J.JDA.2007.01.002.





\item[] 
[19]

A. Dudek, A. Frieze, A. Ruciński, and M. Šileikis.


Embedding the Erdős–Rényi hypergraph into the random
regular hypergraph and hamiltonicity.


\emph{Journal of Combinatorial Theory, Series B}, 122:719–740, 2017.


doi:10.1016/J.JCTB.2016.09.003.





\item[] 
[20]

J. Friedman.


A proof of Alon’s second eigenvalue conjecture and related
problems.


\emph{Memoirs of the American Mathematical Society}, 195(910), 2008.





\item[] 
[21]

P. Gao.


Triangles and subgraph probabilities in random regular graphs.


\emph{Electronic Journal of Combinatorics}, 31(1):P1.2, 2024.


doi:10.37236/10281.





\item[] 
[22]

P. Gao, M. Isaev, and B. D. McKay.


Sandwiching random regular graphs between binomial random graphs.


In \emph{Proceedings of the Fourteenth Annual ACM-SIAM Symposium on
Discrete Algorithms (SODA’20)}, pages 690–701, 2020.





\item[] 
[23]

P. Gao, M. Isaev, and B. D. McKay.


Sandwiching dense random regular graphs between binomial random
graphs.


\emph{Probability Theory and Related Fields}, 184(1–2):115–158,
2022.





\item[] 
[24]

P. Gao, M. Isaev, and B. D. McKay.


Kim-Vu’s sandwich conjecture is true for $d\gg\log^{4}n$.


\emph{arXiv preprint}, 2023.


arXiv:2011.09449.





\item[] 
[25]

J. Gaudio, M. Z. Rácz, and A. Sridhar.


Average-case and smoothed analysis of graph isomorphism.


\emph{Annals of Applied Probability}, 35(2):1373–1406, 2025.





\item[] 
[26]

Y. He.


Spectral gap and edge universality of dense random regular graphs.


\emph{Communications in Mathematical Physics}, 405(181):1–40, 2024.





\item[] 
[27]

H. A. Helfgott, J. Bajpai, and D. Dona.


Graph isomorphisms in quasi-polynomial time.


\emph{arXiv preprint}, 2017.


arXiv:1710.04574.





\item[] 
[28]

N. Immerman and E. S. Lander.


Describing graphs: A first-order approach to graph canonization.


In \emph{Complexity Theory Retrospective}, pages 59–81. Springer, New
York, 1990.





\item[] 
[29]

M. Isaev, T. Makai, B. D. McKay, P. Prałat, J. Tan, and M. Zhukovskii.


Canonical labelling of random regular graphs.


\emph{arXiv preprint}, 2026.


arXiv:2602.17567.





\item[] 
[30]

S. Janson, T. Łuczak, and A. Ruciński.


\emph{Random Graphs}.


Wiley, 2000.





\item[] 
[31]

J. H. Kim, B. Sudakov, and V. Vu.


On the asymmetry of random regular graphs and random graphs.


\emph{Random Structures \& Algorithms}, 21(3–4):216–224, 2002.


doi:10.1002/RSA.10054.





\item[] 
[32]

J. H. Kim, B. Sudakov, and V. Vu.


Small subgraphs of random regular graphs.


\emph{Discrete Mathematics}, 307:1961–1967, 2007.


doi:10.1016/J.DISC.2006.09.032.





\item[] 
[33]

J. H. Kim and V. H. Vu.


Sandwiching random graphs: universality between random graph models.


\emph{Advances in Mathematics}, 188(2):444–469, 2004.





\item[] 
[34]

T. Klimošová, C. Reiher, A. Ruciński, and M. Šileikis.


Sandwiching biregular random graphs.


\emph{Combinatorics, Probability and Computing}, 32(1):1–44, 2023.


doi:10.1017/S0963548322000049.





\item[] 
[35]

M. Krivelevich and B. Sudakov.


Pseudo-random graphs.


In \emph{More Sets, Graphs and Numbers}, volume 15 of \emph{Bolyai
Society Mathematical Studies}, pages 199–262. Springer, 2006.





\item[] 
[36]

M. Krivelevich, B. Sudakov, V. Vu, and N. Wormald.


Random regular graphs of high degree.


\emph{Random Structures \& Algorithms}, 18:346–363, 2001.


doi:10.1002/RSA.1013.





\item[] 
[37]

L. Kučera.


Canonical labeling of regular graphs in linear average time.


In \emph{28th Annual Symposium on Foundations of Computer Science
(SFCS’87)}, pages 271–279, 1987.





\item[] 
[38]

A. Liebenau and N. Wormald.


Asymptotic enumeration of graphs by degree sequence, and the degree
sequence of a random graph.


\emph{Journal of the European Mathematical Society}, 26:1–40, 2024.





\item[] 
[39]

N. Linial and J. Mosheiff.


On the rigidity of sparse random graphs.


\emph{Journal of Graph Theory}, 85(2):466–480, 2017.


doi:10.1002/JGT.22073.





\item[] 
[40]

E. Luks.


Isomorphism of graphs of bounded valence can be tested in polynomial
time.


\emph{Journal of Computer and System Sciences}, 25:42–65, 1982.


doi:10.1016/0022-0000(82)90009-5.





\item[] 
[41]

B. D. McKay.


Subgraphs of random graphs with specified degrees.


\emph{Congressus Numerantium}, 33:213–223, 1981.





\item[] 
[42]

B. D. McKay and N. C. Wormald.


Automorphisms of random graphs with specified degrees.


\emph{Combinatorica}, 4:325–338, 1984.





\item[] 
[43]

B. D. McKay and N. C. Wormald.


Asymptotic enumeration by degree sequence of graphs of high degree.


\emph{European Journal of Combinatorics}, 11:565–580, 1990.


doi:10.1016/S0195-6698(13)80042-X.





\item[] 
[44]

B. D. McKay and N. C. Wormald.


Asymptotic enumeration by degree sequence of graphs with degrees
$o(\sqrt{n})$.


\emph{Combinatorica}, 11:369–382, 1991.





\item[] 
[45]

A. Sarid.


The spectral gap of random regular graphs.


\emph{Random Structures \& Algorithms}, 63(2):557–587, 2023.


doi:10.1002/RSA.21150.





\item[] 
[46]

K. Tikhomirov and P. Youssef.


The spectral gap of dense random regular graphs.


\emph{Annals of Probability}, 41(1):362–419, 2019.





\item[] 
[47]

O. Verbitsky and M. Zhukovskii.


Canonical labeling of sparse random graphs.


In \emph{42nd International Symposium on Theoretical Aspects of
Computer Science (STACS’25)}, pages 75:1–75:20, 2025.





\item[] 
[48]

V. Vu.


Random discrete matrices.


In \emph{Horizon of Combinatorics}, volume 17 of \emph{Bolyai Society
Mathematical Studies}, pages 257–280. Springer, 2008.


doi:10.1007/978-3-540-77200-2.





\item[] 
[49]

B. Weisfeiler and A. Leman.


The reduction of a graph to canonical form and the algebra which
appears therein.


\emph{Nauchno-Technicheskaya Informatsia}, 9(2):12–16, 1968.





\end{itemize}


\end{document}
